% concatenated from main.tex, header.tex, appendix.tex

\documentclass[10pt]{article} % For LaTeX2e
\usepackage[accepted]{tmlr}


\usepackage{amsfonts}    
\usepackage{color}
\usepackage{multirow}
\usepackage{graphicx}
\usepackage{amssymb,amsmath} % NOT COMATIBLE WITH svjour3
\usepackage{enumitem}
\usepackage{float}
%\usepackage{algorithm2e}
%\usepackage{algorithmic}
\sloppy

% REFERENCING
\usepackage[hidelinks]{hyperref}
\usepackage[capitalize, nameinlink]{cleveref}
\crefname{assumption}{Assumption}{Assumptions}


\usepackage{subcaption}
\usepackage{caption}

\usepackage{enumitem}


\newcommand\numberthis{\addtocounter{equation}{1}\tag{\theequation}}

\graphicspath{{figures/}}

\newcommand{\step}{\mbox{H}}
\newcommand{\red}[1]{\textcolor[rgb]{0.635,0.0780,0.1840}{#1}}
\parindent 0ex
\parskip 2ex

\newcommand{\blue}[1]{\textcolor[rgb]{0,0.33,0.56}{#1}}
\newcommand{\green}[1]{\textcolor[rgb]{0,0.65,0.33}{#1}}
\definecolor{shade}{rgb}{0.9,0.9,0.9}
\usepackage{bm}
\renewcommand{\vec}[1]{\bm{#1}}
\newcommand{\mat}[1]{\bm{#1}}
\newcommand{\bx}{\bm{x}}
\newcommand{\by}{\bm{y}}
\newcommand{\bz}{\bm{z}}
\newcommand{\bw}{\bm{w}}
\newcommand{\be}{\bm{e}}
\newcommand{\m}{\bm{m}}
\newcommand{\bg}{\bm{g}}
\newcommand{\bq}{\bm{q}}
\newcommand{\bt}{\bm{\theta}}

%% Robert commands
\newcommand{\dotprod}[1]{\left< #1\right>} 
\DeclareMathOperator{\sgn}{sgn} 
\DeclareMathOperator{\argmininn}{argmin} 
\DeclareMathOperator{\argmaxinn}{argmax} 
\newcommand{\argmin}[1]{ \underset{#1}{\argmininn} \;}
\newcommand{\argmax}[1]{ \underset{#1}{\argmaxinn}\;}
\usepackage[hidelinks]{hyperref}
\newcommand{\cN}{\mathcal{ N}}
\newcommand{\cD}{\mathcal{ D}}
\newcommand{\cP}{\mathcal{P}}
\newcommand{\mH}{\mat{ H}}
\newcommand{\mM}{\mat{ M}}
\newcommand{\mA}{\mat{ A}}
\newcommand{\mX}{\mat{ X}}
\newcommand{\mW}{\mat{ W}}
\newcommand{\mI}{\mat{ I}}
\newcommand{\mS}{\mat{ \Sigma}}
\newcommand{\mD}{\mat{ D}}
\newcommand{\R}{\mathbb{R}}
\DeclareMathOperator{\traceOp}{tr}  
\providecommand{\trace}[1]{\traceOp\left(#1\right)}
\newcommand{\E}[1]{\mathbb{E}\left[#1\right] } 
\newcommand{\EE}[2]{\mathbb{E}_{#1}\left[#2\right] }
\newcommand{\norm}[1]{ \|#1 \|}

% Fabians commands
\newcommand{\prox}[1]{\mathrm{prox}_{#1}}
\newcommand{\iprod}[2]{\langle #1, #2 \rangle}
\newcommand{\N}{\mathbb{N}}
\newcommand{\dist}{\mathcal{D}}
\newcommand{\median}{\texttt{median}}
\DeclareMathOperator*{\mean}{\texttt{mean}}
\newcommand{\Proj}[1]{\mathrm{Proj}_{#1}} % projection
\DeclareMathOperator{\Diag}{Diag} %creates matrix from vector

%switch phi and varphi
\let\temp\phi
\let\phi\varphi
\let\varphi\temp

% Guillaume
\DeclareMathOperator{\SGN}{SGN} 
\DeclareMathOperator{\dom}{dom} 
\DeclareMathOperator{\Dist}{dist} 

% Project specific commans
\newcommand{\agg}[1]{\widehat{\texttt{median}}_{#1}}
\newcommand{\aggt}[1]{\texttt{aggregate}_{#1}}
\newcommand{\samplemedian}[1]{\texttt{median}_{#1,i\in[n]}}
\newcommand{\cclip}[1]{\mathrm{clip}_{#1,1}}
\newcommand{\vclip}[1]{\mathrm{clip}_{#1,2}}
\newcommand{\ClipTO}{\texttt{Clip21}}

% Algorithm names
\newcommand{\SGD}{\texttt{SGD}}
\newcommand{\SGDM}{\texttt{SGD-M}}
\newcommand{\SPP}{\texttt{SPP}}
\newcommand{\Adam}{\texttt{Adam}}
\newcommand{\SMGD}{\texttt{SMGD}}

\newcommand{\VC}{\texttt{VClip}}
\newcommand{\CC}{\texttt{CClip}}
\newcommand{\HC}{\texttt{Huber}}

%
\usepackage{amsthm}
\usepackage{amsbsy}

\usepackage{mathtools}
\usepackage{mdframed} 
\usepackage{thmtools}
\usepackage{thm-restate}
\usepackage{sidecap} % figures with caption on the side
\sidecaptionvpos{figure}{t}


\definecolor{shadecolor}{gray}{0.90}
\declaretheoremstyle[
headfont=\normalfont\bfseries,
notefont=\mdseries, notebraces={(}{)},
bodyfont=\normalfont,
postheadspace=0.5em,
spaceabove=6pt,
mdframed={
  skipabove=8pt,
  skipbelow=8pt,
  hidealllines=true,
  backgroundcolor={shadecolor},
  innerleftmargin=4pt,
  innerrightmargin=4pt}
]{shaded}


\declaretheorem[style=shaded,within=section]{definition}
\declaretheorem[style=shaded,sibling=definition]{theorem}
\declaretheorem[style=shaded,sibling=definition]{proposition}
\declaretheorem[style=shaded,sibling=definition]{assumption}
\declaretheorem[style=shaded,sibling=definition]{corollary}
\declaretheorem[style=shaded,sibling=definition]{conjecture}
\declaretheorem[style=shaded,sibling=definition]{lemma}
\declaretheorem[style=shaded,sibling=definition]{remark}
\declaretheorem[style=shaded,sibling=definition]{example}

% Commenting
\usepackage[colorinlistoftodos]{todonotes}


\definecolor{todored}{RGB}{189, 30, 30}
\newcommand\revise[1]{{\color{todored}{#1}}}

%%%%%%%%%%%%%%%%%%%%%%%%
%%% STYLING
%%%%%%%%%%%%%%%%%%%%%%%%
\definecolor{kleinblue}{RGB}{0, 47, 167}
\hypersetup{
	colorlinks = True,
	linkcolor = kleinblue,
	citecolor=kleinblue,
	urlcolor=kleinblue
} 



\title{Tracking the Median of Gradients with a Stochastic Proximal Point Method
}

% Authors must not appear in the submitted version. They should be hidden
% as long as the tmlr package is used without the [accepted] or [preprint] options.
% Non-anonymous submissions will be rejected without review.

\author{\name Fabian Schaipp \email fabian.schaipp@inria.fr \\
      \addr Technical University of Munich and Inria Paris, ENS, PSL Research University
      \AND
      \name Guillaume Garrigos \email garrigos@lpsm.paris \\
      \addr Université Paris Cité and Sorbonne Université, CNRS \\
      Laboratoire de Probabilités, Statistique et Modélisation, F-75013 Paris, France
      \AND
      \name Umut \c{S}im\c{s}ekli \email umut.simsekli@inria.fr\\
      \addr Inria, CNRS, ENS, PSL Research University \\
      Paris
      \AND
      \name Robert M.\ Gower \email rgower@flatironinstitute.org\\
      \addr CCM, Flatiron Institute,  Simons Foundation \\
      New York City
      }

% The \author macro works with any number of authors. Use \AND 
% to separate the names and addresses of multiple authors.

\newcommand{\fix}{\marginpar{FIX}}
\newcommand{\new}{\marginpar{NEW}}

\def\month{10}  % Insert correct month for camera-ready version
\def\year{2025} % Insert correct year for camera-ready version
\def\openreview{\url{https://openreview.net/forum?id=WMxLEgYGxu}} % Insert correct link to OpenReview for camera-ready version

% Added for footnotes without markers
\newcommand\blfootnote[1]{%
  \begingroup
  \renewcommand\thefootnote{}\footnote{#1}%
  \addtocounter{footnote}{-1}%
  \endgroup
}

\begin{document}


\maketitle

\begin{abstract}
There are several applications of stochastic optimization 
where one can benefit from a robust estimate of the gradient. For example, domains such as distributed learning with corrupted nodes, the presence of large outliers in the training data, learning under privacy constraints, or even heavy-tailed noise due to the dynamics of the algorithm itself. Here we study  \SGD{} with robust gradient estimators based on estimating the median.

We first derive iterative methods based on the  stochastic proximal point method for computing the  median gradient and generalizations thereof. 
Then we propose an algorithm  estimating the median gradient across \emph{iterations}, and find that several well known methods are particular cases of this framework.
For instance, we observe that different forms of clipping allow to compute online estimators of the \emph{median} of gradients, in contrast to (heavy-ball) momentum, which corresponds to an online estimator of the \emph{mean}.
Finally, we provide a theoretical framework for an algorithm computing the median gradient across \emph{samples}, and show that the resulting method can converge even under heavy-tailed, state-dependent noise.
\end{abstract}

\section{Introduction}
  
Many problems in machine learning can be represented by an optimization problem
\begin{equation*} \label{eq:prob}
 \min_{\bw\in \R^d} \ell(\bw),
\end{equation*} 
where $\vec{w}$ is a vector of parameters and $\ell :\mathbb{R}^d \to \mathbb{R}$ is a loss function. 
To tackle this problem, gradient-based optimization algorithms have been the de-facto choice in many application domains \citep{Nemirovsky1983,Bottou2010,Bottou2018}, where we are only given access to a \emph{noisy oracle} of the  gradient %, i.e., 
\begin{equation}\label{eqn:noisy-oracle}
\vec{g}_t = \nabla \ell(\vec{w}_t) + \vec{\xi}_t(\vec{w}_t). 
\end{equation}
Here $t$ denotes the iterations and $\vec{\xi}_t(\vec{w}_t) \in \R^d $ is a noise vector that is sampled from a distribution that depends on the parameter $\vec{w}_t$.
% 
With this noise oracle, stochastic gradient descent (\SGD) with learning rate $\eta_t$ is given by 
\begin{align}\tag{\SGD} \label{eq:sgd}
\bw_{t+1} = \bw_t - \eta_t \bg_t.
\end{align}
While~\ref{eq:sgd} has proven useful in numerous applications, its  performance  heavily relies on the `quality' of the noisy oracles $\vec{g}_t$, and can tragically degrade when the noise $\vec{\xi}_t$ has significant outliers, exhibits heavy tails (that is, $\|\vec{\xi}_t\|$ is very large with non-negligible probability), or has been corrupted in an adversarial way. We provide two examples.


\begin{enumerate}[itemsep=0pt,topsep=0pt,leftmargin=.2in, label=(\roman*)]
\item \textbf{Heavy-tailed gradients.} Recent studies have provided theoretical and empirical evidence that 
    heavy tails can naturally arise in stochastic optimization. From a theoretical perspective \citet{pmlr-v139-gurbuzbalaban21a,hodgkinson2021multiplicative,pavasovic2023approximate} showed that the parameter-dependent nature of $\vec{\xi}_t$ can result in heavy tails, even in very basic settings such as linear regression with Gaussian data. On the other hand, it has been empirically observed that the gradients for training transformer architectures on language tasks are more heavy-tailed compared to, for example, convolutional models for image data \citep{zhang2020adaptive,kunstner2023noise}. While the presence of heavy-tails in gradients might be beneficial in certain settings \citep{simsekli2020hausdorff}, from an optimization perspective, it mostly introduces non-trivial challenges, which might even make the algorithm diverge unless additional care is taken \citep{zhang2020adaptive,Gorbunov2020}.  
\item \textbf{Corrupted nodes.} A well-known problem in distributed learning is 
when some computation nodes are malicious and can communicate adversarial updates, which results in inaccurate gradient oracles that can significantly misguide the optimization procedure. 
Similar to the heavy-tailed setting, depending on 
how malicious the nodes are,
\SGD{} can be impractical unless certain modifications are made~\citep{ICML-2018-MhamdiGR}.
We refer to \citet{Karimireddy2021} and references therein for an overview of more robust aggregation techniques.
\end{enumerate}


In order to make the optimization algorithms more robust to such noisy oracles, many techniques have been developed for different domains under various conditions. In the context of heavy-tailed gradients, clipping is often employed \citep{Zhang2020,Puchkin2023,Koloskova2023}. In distributed learning robust aggregation rules are needed to protect against corrupted or malicious nodes \citep{Data2018,Karimireddy2021,khirirat2023clip21}. This line of research is also closely related to error feedback in federated learning \citep{seide2014-bit,Karimireddy2019,richtarik2021ef} and learning under differential privacy constraints \citep{khirirat2023clip21}.


\paragraph{Robustness of median.} While attracting increasing attention in stochastic optimization thanks to modern machine learning applications, taming the impact of strong noise has already been considered in various other domains. Indeed, estimation under heavy-tails and corruptions has been a long-standing topic in robust statistics \citep{Huber1981}. Being a vast research field that has produced numerous algorithms designed for different tasks, one of the recurring themes in robust statistics is the use of the \emph{sample median} and its  variations as opposed to the \emph{sample mean}~\citep{Minsker2015}, whenever an aggregation of random variables is needed. For instance, in the presence of heavy tails, the sample median has been shown to be a significantly more robust estimator of the true mean, whereas the sample mean can be vulnerable to strong noise \citep{Lugosi2019}.

\paragraph{Sample median \SGD{}.} Inspired by tools from robust statistics, some notions of median have been utilized in stochastic optimization as well~\citep{pmlr-v80-yin18a,Alistarh2018,pmlr-v151-acharya22a}.
One such approach is to compute the median over a finite sample of gradients. 
While this approach has paved the way for powerful optimization algorithms in terms of robustness, it often introduces a significant computational burden, since a multivariate median needs to be computed at every iteration \citep{Weiszfeld2009,Vardi2000}.


\textbf{Motivation.} Clipping-based and error feedback approaches on the one hand, and median-based approaches on the other hand are motivated through different mathematical frameworks and appear to use different algorithmic tools. In this study, our main goal is to bridge the gap between these two seemingly different branches and to bring a unified theoretical perspective that can shed further light on both directions. The main tool to build this bridge between stochastic optimization and robust statistics will be the Stochastic Proximal Point method (\SPP) \citep{Asi2019,Davis2019,Toulis2020} for solving a later-specified gradient estimation problem.   

\textbf{Contributions.} 
%
We first show how heavy-ball momentum \citep{Polyak1964} can be seen as online estimator of the mean gradient using \SPP{}. Here we use the term \emph{online} in the sense that the method is given a  different loss function at each step.


Encouraged by this result, we then focus on online estimation of the \emph{median} gradient, using the same technique. Doing so, we recover several known clipping-based optimization algorithms from distributed learning as special cases, such as~\ClipTO{}~\citep{khirirat2023clip21}, and a variant of \emph{centered clipping} \citep{Karimireddy2021}. We also shown connections to sign-based gradient methods \citep{Riedmiller1993}.

Our proposed framework allows to unify the different motivations for previously developed techniques, and illustrates that clipping-based optimization algorithms essentially run a loose estimation of the median across iterations. In particular, the observation that \emph{centered clipping} and \ClipTO{} are iterative estimators of a geometric median appears to be new, and might have consequences for its application in distributed learning.

We theoretically analyze these online gradient estimation methods in a simplified setting: when allowing multiple gradient samples per iteration, we show that \SGD{} with (approximately computed) median gradients indeed yields a very powerful algorithm, which converges with $1/\sqrt{T}$ rate ($T$ being the number of iterations) even when the gradient oracles are heavy-tailed with diverging second-order moments. We further illustrate that, in addition to having infinite variance, the noise vectors $\vec{\xi}_t(\vec{w})$ can even have a multiplicative dependence on $\vec{w}$, a scenario which is highly challenging and cannot be directly covered by existing theoretical results, see e.g., \citet{wang2021convergence,zhang2020adaptive,Puchkin2023}. 

We finally illustrate our theory on synthetic least-squares experiments where we compare the effectiveness of sample median, sample mean, and several online median estimators. Our results underline that for heavy-tailed noise, using the sample median is highly effective in contrast to the sample mean which is unstable and often does not converge.
Our experiments also show that our online median estimates are robust, 
and require only a single sample per iteration, making it a less expensive alternative to the sample median.
We further compare different clipping techniques for training transformer architectures on language modeling tasks, and show that they can improve upon the performance of \SGD{} with momentum, however the gap is relatively small. 


\textbf{Notation.} Throughout the paper, $\|\cdot\|_p$ denotes the  $\ell_p$-norm, given by $\norm{\vec{z}}_p : = \left( \sum_{i=1}^d|\vec{z}|_i^p \right)^{\frac{1}{p}}$ when $p\in [1,\infty)$, and $\Vert \vec{z} \Vert_\infty = \max\limits_{i=1, \dots, d} \vert \vec{z}_i \vert$ when $p=\infty$. 
When $p=2$ we recover the usual Euclidean $\ell_2$-norm $\Vert \cdot \Vert_2$, which we will often simply denote by $\Vert \cdot \Vert$. For a matrix $\vec{A}$, we denote its Frobenius norm by $\|\vec{A}\|_F := \sqrt{\sum_{i,j}\vec{A}_{ij}^2}$. {For $n \in \mathbb{N}_{+}$, we denote the set $[n]:=\{1,\dots, n\}$.} 




\section{Preliminaries}
%
\subsection{Heavy-tailed Random Variables}

A random variable $X$ is said to be \emph{heavy-tailed} if
the tails of its distribution decay slower
than the tails of an exponential distribution.
One important example is the family of symmetric $\alpha$-stable distributions $\mathcal{P}_{\alpha, \sigma}$ \citep[Def.\ 1.4]{Nolan2020}, parameterized by a stability index $\alpha \in (0, 2]$ and a scale $\sigma > 0$. When $\alpha = 1$, this is the Cauchy distribution; when $\alpha = 2$, it reduces to a Gaussian.
%
The parameter $\sigma$ controls the spread of the distribution. When $\sigma = 1$, we call the distribution \emph{standard} and denote it by $\mathcal{P}_{\alpha} := \mathcal{P}_{\alpha,1}$.
The parameter $\alpha$ on the other hand controls the heaviness of the tails: for $\alpha<2$ the random variable $X\sim \mathcal{P}_{\alpha}$ has infinite variance, whereas for $\alpha\leq 1$ even the mean of $X$ is infinite. However, for any $\alpha$, the median is equal to $0$. 


\subsection{Mean Estimation and (Sample) Median}

Estimating the mean of a distribution, when given a finite number of samples, is one of the core problems studied in statistics. 
If the distribution in question has heavy-tails, or has significant outliers, the sample mean turns out to be a poor estimator. \citet{Lugosi2019} give a survey of alternative and more robust techniques: one such alternative is the median, for which we will introduce basic definitions next. 

Here, we denote by $z$ a real-value random value, and the boldface  $\vec{z}\in \R^d$ a random vector. We always denote the expectation with respect to the distribution of $z$ (or $\vec{z}$ respectively) by $\mathbb{E}$. 

\paragraph{Multivariate median.} The one-dimensional median is defined by
\begin{align}\label{eq:def-one-dim-median}
    \median{}(z) \in 
    \argmin{m \in \mathbb{R}} \E{|m-z|}.
\end{align}
In contrast to the mean, the median is always defined (though it may not be unique) \citep{Cramer2016}.
Since we are interested in approximating the median gradient, 
we need a notion of median that generalizes to random vectors $\vec{z}\in \R^d$. For this, we can define the $\ell_p$-median: \begin{equation}\label{eq:lp-median}
    \median_{\ell_p}(\vec{z}) \in \argmin{\vec{m} \in \R^d} \E{ \norm{\vec{m} - \vec{z}}_p},
\end{equation}
where $\|\vec{m}\|_p =\left( \sum_{j=1}^d \vec{m}_j^p \right)^{\frac{1}{p}}$ is the $\ell_p$ norm.
This recovers several existing notions of a multivariate median: for $p=1$ we obtain a \emph{componentwise median} (or $\ell_1$-median) and for $p=2$ we obtain the \emph{geometric median} (or $\ell_2$-median)  \citep{Frechet1948,Weiszfeld2009,Minsker2015,Cohen2016}. For $d=1$ and every $p \in [1, \infty)$, the $\ell_p$-median coincides with the one-dimensional median in \eqref{eq:def-one-dim-median}.

\paragraph{Sample median.} Given $n$ realizations $\vec{z}^{(1)},\ldots,\vec{z}^{(n)}$ of $\vec{z}$, the sample median is defined as 
\begin{equation}\label{eq:sample-median}
    \samplemedian{\ell_p}(\vec{z}^{(i)}) \in \argmin{\vec{m} \in \R^d} \frac{1}{n}\sum_{i=1}^n \norm{\vec{m} - \vec{z}^{(i)}}_p.
\end{equation}
Clearly \eqref{eq:sample-median} is a special case of \eqref{eq:lp-median}, by setting the distribution of $\vec{z}$ as the empirical measure $\frac{1}{n}\sum_{i=1}^n \delta_{\vec{z}_i}$.

It is generally accepted that the sample median is a more robust estimator than the sample mean \citep{Huber1981,Hampel2005}.
This is nicely illustrated by the notion of a breakdown point \citep{Donoho1983,Lopuhaa1991},  which is the smallest fraction of a sample that, if arbitrarily corrupted, can arbitrarily change the value of the estimator. The breakdown point of the sample median is roughly $\tfrac12$, meaning that at least half of the samples need to be outliers for the median to diverge. The sample mean however can be arbitrarily corrupted by a single outlier, thus its breakdown point is $\tfrac{1}{n}$. 

If the samples $\vec{z}^{(i)}$ are obtained by taking block-wise means over observations, \eqref{eq:sample-median} is called the median-of-means estimator \citep{Nemirovsky1983,Lugosi2019}, which is known to be a robust mean estimator \citep{Lugosi2019}. 


\paragraph{Relation to M-estimation.} More generally, for a function $\dist: \R^d \to \R$ we will consider the problem
\begin{align}\label{prob:gradient-learning}
   \argmin{\vec{m} \in \R^d} \E{ \dist(\vec{m}-\vec{z})}.
\end{align}
We assume that  $\dist$ is a closed,  proper, convex function and that~\eqref{prob:gradient-learning} admits a solution (see \cref{sec:app-more-update-formulas} for a discussion).
For $\dist = \tfrac12\norm{\cdot}_2^2$, the solution to \eqref{prob:gradient-learning} is the mean $\E{\vec{z}}$ (if it exists), and for $\dist=\|\cdot\|_2$ the solution is the geometric median as defined in \eqref{eq:lp-median}. 
Problem~\eqref{prob:gradient-learning}  can be seen as a special case of \emph{M-estimators} \citep[Sec.\ 3.2]{Huber1981}.

The choice of $\dist$ determines properties of the associated estimator:
if the function $\dist$ heavily penalizes large values, like the quadratic  $\dist = \tfrac12\norm{\cdot}_2^2$, then the solution to \eqref{prob:gradient-learning} will be sensitive to outliers. In contrast  $\dist  = \norm{\cdot}_2$ and $\dist  = \norm{\cdot}_1$ grow only linearly, and thus~\eqref{prob:gradient-learning} will be less sensitive to outliers. We return to problem \eqref{prob:gradient-learning} later in \cref{sec:grad-estimation-with-spp}.

\section{Robust Gradient Estimation with Stochastic Proximal Point}
\subsection{Warmup: Momentum as Online Mean Gradient Estimation}\label{sec:momentum-mean-estimation}

When dealing with noisy gradient oracles, a standard technique is to apply (heavy-ball) momentum \citep{Polyak1964}. For momentum coefficient $\beta\in [0,1)$, it is given by
\begin{align}\label{eqn:momentum}\tag{\texttt{SGD-M}}
    \begin{split}
        \vec{m}_{t+1} =  \beta \vec{m}_t + (1-\beta) \vec{g}_t, \\
    \vec{w}_{t+1} = \vec{w}_t - \eta_t \vec{m}_{t+1}.
    \end{split}
\end{align}
Momentum has been empirically shown to improve the performance of \SGD{} in machine learning \citep{sutskever13}, and can be seen as a variance reduction technique \citep{Gower2020}.

It is easy to see that momentum performs an exponentially weighted average over past gradients. In this section, we will derive a new perspective on \ref{eqn:momentum}, namely being an \emph{online estimator of the mean gradient} at a given point. To see this, we recall that for a random variable $\vec{z}\in \R^d$, its mean (if it exists) is the solution to \eqref{prob:gradient-learning} with $\dist = \frac12\|\cdot\|^2$.
The random variable of interest in this context are the noisy gradients $\vec{g}_t$.
We could in principle solve problem \eqref{prob:gradient-learning} with iterative stochastic methods. For this purpose, we consider the stochastic proximal point (\SPP{}) method \citep{Asi2019,Davis2019}. We give a detailed introduction to \SPP{} in the next section.  

For the purpose of this section, it is sufficient to know that the iterates $(\vec{m}_t)$ of \SPP{} with step-size $\tau > 0$ applied to~\eqref{prob:gradient-learning} are given by (details in Appendix \ref{SS:SPP:general D formulas})
\begin{align}\label{eqn:spp-step}
    \vec{m}_{t+1} \; = \; \argmin{\vec{m}} \dist(\vec{m}-\vec{g}_t) +\frac{1}{2\tau}\|\vec{m}-\vec{m}_t\|^2 
 = \vec{g}_t + \prox{\tau \dist}(\vec{m}_t -\vec{g}_t), 
\end{align}
where we denote by $\prox{\dist}$ the proximal operator of $\mathcal{D}$,
\[\prox{ \dist}(\vec{x})  : = \argmin{\by\in \R^d} \dist(\by) +\frac{1}{2} \|\by-\vec{x}\|^2. \]%
% 
For $\dist = \frac12\|\cdot\|^2$, it is easy to compute $\prox{\tau \dist}(\vec{x}) = \frac{1}{1+\tau}\vec{x}$. Thus, one iteration of \SPP{} \eqref{eqn:spp-step} is equal to 
\begin{align}\label{eqn:momentum-as-spp}
    \vec{m}_{t+1} \; = \; \vec{g}_t + \prox{\tau \dist}(\vec{m}_t -\vec{g}_t) = \vec{g}_t + \tfrac{1}{1+\tau}(\vec{m}_t -\vec{g}_t) = (1-\tfrac{\tau}{1+\tau})\vec{m}_t + \tfrac{\tau}{1+\tau} \vec{g}_t.
\end{align}
This is exactly the momentum step in \eqref{eqn:momentum} with $\beta= 1-\tfrac{\tau}{1+\tau}$.
However, in \ref{eqn:momentum} the parameters $\vec{w}_t$ are updated each step as well, and thereby the distribution of the gradients $\vec{g}_t$ changes each step. In conclusion, the momentum operation can be understood as an \emph{online} estimation of the mean gradient.

We end this section by justifying our choice to specifically consider \SPP{} for solving \eqref{prob:gradient-learning}, instead of \SGD{} or other stochastic methods. 
If we were to apply \SGD{} with step size to problem \eqref{prob:gradient-learning}, we would obtain 
\begin{align*}
    \vec{m}_{t+1} \; = (1-\tau) \vec{m}_t + \tau \vec{g}_t.
\end{align*}
This also resembles momentum, but note that the momentum coefficient $1-\tau$ could be negative if $\tau > 1$. 
For \SPP{} on the other hand, the momentum coefficient is correctly parametrized since $1-\tfrac{\tau}{1+\tau}$ is contained in $[0,1)$ for every $\tau > 0$.

In conclusion, we can understand heavy-ball momentum (\ref{eqn:momentum}) as applying \SPP{} to the \emph{mean} estimation problem in an \emph{online} fashion. This opens the question what is the \SPP{} update for estimators that are more robust to heavy-tailed noise, such as the median (or more generally for other choices for $\dist$ that are of interest).
We answer this in the next section.
%
%
\subsection{Clipping as Online Median Gradient Estimation}\label{sec:grad-estimation-with-spp}

We now focus on the following gradient estimation problem:
let $\mathcal{G}$ denote the (fixed) distribution over noisy gradients $\vec{g}$. For $\dist:\R^d \to \R_{\geq 0}$ being a closed, proper, convex function, we consider the problem of approximating \begin{align}\label{prob:gradient-learning-3}
    \argmin{\vec{m} \in \R^d} \EE{\vec{g}\sim \mathcal{G}} 
   { \dist(\vec{m}-\vec{g})}.
\end{align}

We consider stochastic iterative methods for solving~\eqref{prob:gradient-learning-3}, with a special focus on the stochastic proximal point method (\SPP{}). We recall the \SPP{} update being
\begin{align}\label{eqn:online-spp-update}
    \begin{split}
    \vec{m}_{t+1} \; &= \; \argmin{\vec{m}} \dist(\vec{m}-\vec{g}_t) +\frac{1}{2\tau}\|\vec{m}-\vec{m}_t\|^2 \\
    &= \vec{g}_t + \prox{\tau \dist}(\vec{m}_t -\vec{g}_t).
    \end{split}
\end{align}
The \SPP{} update can be seen as an implicit version of \SGD{}, which is hard to compute in general. In cases where \SPP{} has a closed form update, it is often preferred to \SGD{} since is easier to tune~\citep{Asi2019,Milzarek2024}.
Fortunately, as we show next, the \SPP{} method enjoys closed-form updates for all choices of $\dist$ which are of interest in our context, in particular for $\dist \in \{\|\cdot\|_1,\|\cdot\|_2,\tfrac12\|\cdot\|^2\}$. All the details are deferred to \cref{app:sec:SPP-updates-particular-cases}.

\paragraph{Vectorwise clipping.} 
If we choose $\dist = \|\cdot\|_2$, the \SPP{} update applies clipping to the difference between sampled gradient $\vec{g}_t$ and previous estimate $\vec{m}_t$. 

\begin{restatable}{corollary}{corLtwonorm}\label{cor:L2}
For $\dist = \|\cdot\|_2$ update~\eqref{eqn:online-spp-update} is given by 
\begin{align}
        \vec{m}_{t+1} &= \vec{m}_t + \vclip{\tau}(\vec{g}_t - \vec{m}_t),  \label{eq:vclip-update}
        % \\ &= 
\end{align}
where $\vclip{\tau}(\vec{v}) :=  \frac{\tau }{\max\{\tau, \|\vec{v}\|_2 \}} \vec{v}$.
\end{restatable}
We can rewrite update \eqref{eq:vclip-update} as $\vec{m}_{t+1} = \beta_t \vec{m}_t + (1-\beta_t) \vec{g}_t$ with $\beta_t:= 1-\tfrac{\tau}{\max\{\tau, \|\vec{g}_t-\vec{m}_t\|_2\}}$. In other words, online estimation of the geometric median ($\dist = \|\cdot\|_2$) with \SPP{} is equal to a cautious version of momentum, as samples which are too far from $\vec{m}_t$ will
be down weighted. This allows the iterates to be more robust to outliers.

\paragraph{Componentwise clipping.}  If we choose $\dist= \|\cdot\|_1$, \SPP{} performs componentwise clipping instead. Here as well, clipping protects against large individual entries by shrinking them independently.
\begin{restatable}{corollary}{corLonenorm}\label{cor:L1}
For $\dist= \|\cdot\|_1$ update~\eqref{eqn:online-spp-update} is given by
    \begin{align}
        \vec{m}_{t+1} = \vec{m}_t + \cclip{\tau}(\vec{g}_t - \vec{m}_t),  \label{eq:cclip-update}
    \end{align}
     where $\cclip{\tau}(\vec{v}) := \left(  \min\{\max\{\vec{v}_i, - \tau\} , \tau\}\right)_{i=1}^d$.
\end{restatable}

\paragraph{Huber function.} For update \eqref{eq:vclip-update}, it is possible that $\vec{m}_{t+1} = \vec{g}_t$, thus no momentum. To allow for a soft transition, we can choose $\dist$ to be the Huber function. For $\mu > 0$, it is given by
\begin{align*}
    H_{\mu}: \R^d \to \R, \quad H_{\mu}(\vec{z}) = \begin{cases}
    \frac{1}{2}\|\vec{z}\|^2 &\quad \|\vec{z}\| \leq \mu,\\
    \mu \|\vec{z}\| - \frac{\mu^2}{2} &\quad \text{else}.
\end{cases}
\end{align*}
The Huber function operates like the squared loss for arguments with small norm, and like the $\ell_2$-norm for arguments with large norms (potentially outliers). 

\begin{restatable}{corollary}{corLhuber}\label{cor:huber}
For $\dist = H_{\mu}$, update~\eqref{eqn:online-spp-update} is given by
\begin{align} \label{eq:huber-update}
    \vec{m}_{t+1} &= \vec{g}_t + \beta_t (\vec{m}_t -\vec{g}_t) 
    =    \beta_t \vec{m}_t 
    + (1-\beta_t) \vec{g}_t,
\end{align}
where $\beta_t := 1-\tfrac{\mu\tau}{\max\{\|\vec{m}_t- \vec{g}_t\|,~ \mu(1+\tau)\}}$.
\end{restatable}



In update \eqref{eq:huber-update} the coefficient $\beta_t$ is always positive. Thus, $\dist = H_{\mu}$ is a tradeoff between momentum with fixed coefficient ($\dist=\frac12\|\cdot\|^2$) and clipping ($\dist = \|\cdot\|_2$). However, it comes at the cost of choosing/tuning an additional hyperparameter $\mu$. In \cref{sec:SGDmedian}, we additionally compare the updates when using stochastic subgradient descent instead of \SPP{} for solving \eqref{prob:gradient-learning-3}, showing that sign-\SGD{} is related to $\ell_1$-median estimation.

\section{Gradient Estimator in the Wild}\label{sec:grad-est-in-the-wild}

Our final objective is to use robust gradient estimators within a \SGD-type method. Because at each iteration the weights $\bw_t$ are updated, the distribution of the gradients (denoted by $\mathcal{G}$ in the previous section) also changes at each iteration of \SGD{}. 
Our strategy is  to interweave updates in the weights $\bw_t$ with \SPP{} updates in the gradient estimators $\vec{m}_t.$ 
%That is, we consider methods of the form
That is, we sample \emph{one} stochastic gradient $\bg_t$ per iteration  and set
\begin{align} \label{eq:online-update}
\begin{array}{llr}    
    \vec{m}_{t+1} &=\quad 
    \vec{g}_t + \prox{\tau\dist}(\vec{m}_t-\vec{g}_t), \\%& \qquad \mbox{\ttfamily\textcolor{gray}{\# Update gradient estimate}} \\
    \vec{w}_{t+1} &= \quad\vec{w}_t - \eta_t \vec{m}_{t+1}. \\%& \mbox{ \ttfamily\textcolor{gray}{\# Update parameters}}  
    \end{array}
\end{align} 

With the results from \cref{sec:momentum-mean-estimation,sec:grad-estimation-with-spp}, we immediately obtain the following special cases: 

(i) \quad if $\dist = \frac{1}{2}\|\cdot\|^2$ : \
     $\begin{cases}
         \vec{m}_{t+1}  =  (1-\tfrac{\tau}{1+\tau}) \vec{m}_t + \tfrac{\tau}{1+\tau} \vec{g}_t \\
         \bw_{t+1}  = \bw_t - \eta_t \m_{t+1},
     \end{cases}$
     
(ii) \quad if $\dist = \Vert \cdot \Vert_p$, $p \in \{1,2 \}$  :   \
     $\begin{cases}
         \vec{m}_{t+1} = \vec{m}_t + {\rm clip}_{p,\tau}(\vec{g}_t - \vec{m}_t) \\
         \bw_{t+1}  = \bw_t - \eta_t \m_{t+1}.
     \end{cases} $ 

Equation (i) is \SGD{} with momentum where the \SPP{} step size $\tau$ defines the momentum coefficient.
On the other hand, (ii) corresponds to clipping, and here the \SPP{} step size determines the clipping threshold.
In other words, \SGD{} with momentum is to online mean estimation what clipping is to online median estimation.

A slightly different approach would be to restart $\m_t=\vec{0}$ in every iteration (\emph{cold start}). 
In this case, (i) reduces to \SGD{} with no momentum, while in (ii) the clipping operator is applied to $\vec{g}_t$, instead of $\vec{g}_t-\vec{m}_t$, which is also a common technique in practice \citep{Zhang2020}. We discuss this relation below in more detail.

\paragraph{Connection to Error Feedback.}
Methods such as~\eqref{eq:online-update} are also used in the distributed learning setting to improve communication and protect against corrupted nodes~\citep{Karimireddy2021}. These methods are often called  Error Feedback (\texttt{EF}) methods~\citep{seide2014-bit,richtarik2021ef}
and generally involve an update of the form
$\vec{m}_{t+1} = \vec{m}_t + \mathcal{C}_t (\vec{g}_t-\vec{m}_t)$
where $\mathcal{C}_t :\R^d \mapsto \R^d$ is a compression operator.
Projections onto balls can be seen as a compression, since one can encode the resulting vector in fewer bits~\citep{Safaryan2021}.
In particular, centered clipping by \citet{Karimireddy2021}, and the \texttt{Clip21} method~\citep{khirirat2023clip21} are using variations of update~\eqref{eq:vclip-update} in the distributed setting, and are thus `secretly estimating' the geometric median of the gradients returned across all nodes.

\paragraph{\SGD{} with standard clipping.}
Even in the non-distributed setting, a standard technique to avoid training instabilities is to directly clip the gradient $\vec{g}_t$ \citep{Zhang2020}. For momentum coefficient $\beta \in [0,1)$ and $c>0$, let 
\begin{align}\label{eqn:clipped-sgdm}
    \begin{split}
    \vec{m}_{t+1} &= \beta \vec{m}_t + (1-\beta) \min\{1,\tfrac{c}{\|\vec{g}_t\|}\}\vec{g}_t, \\
    \vec{w}_{t+1} &= \vec{w}_t - \eta_t \vec{m}_{t+1}.
    \end{split}
\end{align}
We will refer to \eqref{eqn:clipped-sgdm} as \texttt{clipped-SGD}.
The method has been studied extensively, often for the case $\beta=0$, see \citet{Zhang2020,Gorbunov2020,Koloskova2023,Puchkin2023} and references therein.

As explained above, when $\beta=0$, then \eqref{eqn:clipped-sgdm} is the same as \eqref{eq:vclip-update} with cold start (resetting $\vec{m}_t=\vec{0}$ in every iteration); but this does not allow for momentum, which is usually used in practice on top of clipping. 
However, we can derive \texttt{clipped-SGD} \eqref{eqn:clipped-sgdm} also as online \SPP{} for computing a specific robust mean estimator, known as the \textit{Catoni-Giulini mean estimator} \citep{Catoni2018,Lugosi2019}. For this, consider the problem
\begin{align}\label{prob:clipped-sgd}
    \argmin{\vec{m} \in \R^d} \EE{\vec{g}\sim \mathcal{G}}{\tfrac12 \|\vec{m} - \min\{1, \tfrac{c}{\|\vec{g}\|}\}\vec{g}\|^2}.
\end{align}
The solution to \eqref{prob:clipped-sgd} is $\EE{\vec{g}\sim \mathcal{G}}{\min\{1, \tfrac{c}{\|\vec{g}\|}\}\vec{g}}$, which is exactly the Catoni-Giulini mean estimator.
Now, one step of \SPP{}~\eqref{eqn:online-spp-update} applied to~\eqref{prob:clipped-sgd} gives \eqref{eqn:clipped-sgdm} with $\beta=\tfrac{1}{1+\tau}$.

\paragraph{Outline of remaining paper.} Before comparing these different online gradient estimation methods in experiments, we will discuss convergence guarantees: ideally, we would like to analyze the online method \eqref{eq:online-update}. However, when $\dist$ is not strongly convex this analysis turns out to be difficult, mainly due to the fact that the distribution of $\vec{g}_t$ changes every step. Hence, the next section will analyze the online median estimator in a simplified setting, where we have access to multiple samples each iteration, and can (approximately) compute their sample median.
We focus on settings where the noise in $\vec{g}_t$ is heavy-tailed, and prove that in such scenarios the sample median can still guarantee convergence (while the sample mean fails).

Finally, our experiments in \cref{sec:experiments} will show that the online versions of the sample median method (\VC{} and \CC{}) show superior performance as well in the heavy-tailed scenario, justifying the online estimation. 


\section{Theory for the Sample Median Gradient Method}\label{sec:sample-median-methods}

In this section, we showcase that \SGD{} with an (approximate) sample median is robust to heavy-tailed noise, whereas \SGD{} using the sample mean can fail.
Consider the following setup: at every iteration $t \geq 0$, we have access to $n$ noisy oracle gradients $\big(\vec{g}_t^{(1)},\ldots,\vec{g}_t^{(n)}\big)$, where
\[  \vec{g}_t^{(i)} = \nabla \ell(\bw_t) + \vec{\xi}^{(i)}_t(\vec{w}_t) \quad \mbox{for} \quad  i=1,\dots, n.\] 
Here $\{\vec{\xi}^{(i)}_t(\vec{w}_t) \in \R^d \}_{i=1}^n$ is a set of  i.i.d. noise vectors. Given these sampled gradients, the typical approach of \ref{eq:sgd} would be to use their sample mean (average) $\vec{g}_t = \frac1{n} \sum_{i=1}^n \vec{g}_t^{(i)}$ as update direction. However, \citet[Remark 1]{zhang2020adaptive} show that this fails to converge for any step size under heavy-tailed noise.

Alternatively, to improve robustness, we analyze the update
\begin{align}\label{eqn:approx-median-descent}
    \vec{w}_{t+1} = \vec{w}_t - \eta (\vec{m}_t + \vec{e}_t), 
\end{align}
where $\vec{m}_t=\samplemedian{\ell_p}(\vec{g}_t^{(i)})$, and $\vec{e}_t\in \R^d$ represents a random error we may commit in computing the sample median. We refer to method \eqref{eqn:approx-median-descent} as \emph{Sample Median Gradient Descent} (\SMGD{}).

\vspace{1ex}

\begin{assumption}
\label{ass:uniform-bias-variance-bound-median}
    There exists $\sigma_1,\sigma_2,\delta_1, \delta_2\geq0$ such that, for every $\bw \in \mathbb{R}^d$, 
      the sample median $\m = \samplemedian{\ell_p}(\vec{g}^{(i)})$ verifies 
\begin{align}\label{ass:noise-approx-median}
    \begin{split}
    &\Vert \E{\vec{m}} - \nabla \ell(\vec{w}) \Vert^2 \leq \delta_1^2 + \delta_2 \|\nabla \ell(\bw)\|^2, \\
    &\E{ \Vert \vec{m} - \E{\vec{m}} \Vert^2  } \leq \sigma_1^2 + \sigma_2 \|\nabla \ell(\bw)\|^2.
    \end{split}
\end{align}
\end{assumption}

\begin{assumption}
\label{Ass:errors bounded expectation}
    There exists $\varepsilon > 0$ such that, for every $t \geq 0$, the error in computing the sample median has finite variance: 
    $\mathbb{E}\left[  \Vert \vec{e}_t \Vert^2 \ | \ \bg_t^{(1)}, \dots, \bg_t^{(n)} \right] \leq \varepsilon^2$.
\end{assumption}

Based on these two assumptions, we establish the following complexity results for~\eqref{eqn:approx-median-descent}.
\begin{restatable}{proposition}{propapproxmed}
\label{prop:approx-median-descent}
Let $\ell : \mathbb{R}^d \to \mathbb{R}$ be $L$-smooth and let \cref{ass:uniform-bias-variance-bound-median} hold with
$\delta_2 \leq \frac18$.
Take $\eta \leq \frac{1}{8L(1+2\sigma_2+2\delta_2)}$ and consider $\bw_t$ a sequence generated by  \eqref{eqn:approx-median-descent} where $\vec{e}_t$ verifies \cref{Ass:errors bounded expectation}.
Then, for every $T \geq 1$ it holds
\begin{equation*}
    \min_{0\leq t\leq T-1} \E{\|\nabla \ell(\bw_t)\|^2}
    =
    \mathcal{O}\left( \tfrac{1}{\eta T} + \eta + \delta_1^2 + \varepsilon^2 \right).
\end{equation*}
\end{restatable}
%

\begin{restatable}{proposition}{CVSMGDconvexlipschitz}\label{P:CV SMGD convex lipschitz}
    Let $\ell : \mathbb{R}^d \to \mathbb{R}$ be convex and $G$-Lipschitz, with ${\rm{argmin}}~\ell \neq \emptyset$. 
    Let \cref{ass:uniform-bias-variance-bound-median,Ass:errors bounded expectation} hold.
    Without loss of generality assume that $\delta_2=\sigma_2=0$.
    Then, for $\eta = \frac{1}{(\delta_1+\varepsilon)T + \sqrt{T}}$ and every $T \geq 1,$ the  iterates $\bw_t$ from \eqref{eqn:approx-median-descent} satisfy
    \begin{equation*}
        \EE{}{\ell(\bar \bw_T) - \inf \ell}
        =
        \mathcal{O}\left( \frac{1}{\sqrt{T}} + \delta_1 + \varepsilon \right),
    \end{equation*}
    where $\bar \bw_T := \tfrac1T \sum_{t=0}^{T-1} \theta^t \bw_t$ with $\theta = \tfrac{T}{T+2}$. 
\end{restatable}

The proofs follow techniques from the analysis of biased \SGD{} (see e.g.\ \citet{DemMalSokRic23}) and are provided in \cref{sec:missing-proofs}. 
Let us now discuss our main \cref{ass:uniform-bias-variance-bound-median}. Assume for simplicity that $\delta_2=\sigma_2=0$ (for instance if $\ell$ is Lipschitz), that $\vec{e}_t = 0$ (the sample median is computed exactly), and consider that $\vec{g}_t^{(i)} = \nabla \ell(\vec{w}_t) + \vec{\xi}_t^{(i)}$, where $\vec{\xi}_t^{(i)}$ is a noise term. 
Then, \cref{ass:uniform-bias-variance-bound-median}
becomes\footnote{See  \cref{lem:noise-assum-equiv} in the appendix for a proof.} 
\begin{align}\label{eqn:sample-median-noise-cond}
    \E{\|\samplemedian{\ell_p}(\vec{\xi}_t^{(i)})\|^2} \leq \nu^2 < + \infty.
\end{align}
In other words, we need the second moment of the sample median of $\vec{\xi}_t^{(i)}$ to be uniformly bounded.
We want to stress that our choice of using the median is crucial when the $\vec{\xi}_t^{(i)}$ come from a heavy-tailed distribution, such as a (univariate) standard $\alpha$-stable distribution, with $\alpha\in (1,2)$.
Indeed, in this case we know that the variance of the sample median is \emph{finite} for sufficiently large sample size $n$ (see \citet[Theorem 2.2]{Bickel1967} and \citet[Theorem 1.2]{Nolan2020})\footnote{We know more:
The sample median of a parent distribution with median $\xi$ is asymptotically normal with location $\xi$ \citep{Chu1955}. For  Cauchy distributions ($\alpha=1$), the sample median has a finite variance for $n\geq 5$ \citep[Remark 1]{Chu1955}, while first and second moment of the sample mean are undefined.}.
On the contrary, taking $\vec{m}_t$ as the usual sample \emph{mean} causes trouble, as the variance of the sample mean is \emph{infinite}\footnote{This is due to the fact that, by definition of $\alpha$-stability, the sample mean of i.i.d sample follows the same $\alpha$-stable distribution (up to affine translation).} everywhere. 
See  \cref{fig:moment-sample-median-mean} for an illustration of this phenomenon.


This simple discussion illustrates that \cref{ass:uniform-bias-variance-bound-median} can be satisfied under mild conditions on the noise. 
We exploit this in the next corollary, which shows that the $\ell_1$-sample median can even accommodate state-dependent noise. 
Let us define first our model for such noise.

\begin{restatable}{assumption}{Assstablenoise}
\label{Ass:oracle with linear stable noise}
    For every $\bw \in \mathbb{R}^d$, 
    we have that
    $\bg = \nabla \ell(\vec{w}) + \vec{\Sigma}(\bw)\vec{\zeta}$, 
    where $\vec{\Sigma}(\bw) \in \mathbb{R}^{d \times d}$ is a matrix and $\vec{\zeta}$ is a random vector whose components are drawn i.i.d.\ from a standard symmetric $\alpha$-stable distribution with $\alpha~>~1$.
    There exists constants $c_1, c_2 >0$ such that
    $$\Vert \vec{\Sigma}(\vec{w})\Vert_F^2 \leq c_1 + c_2 \|\nabla \ell(\vec{w})\|^2 \quad \mbox{for all  } \vec{w} \in \R^d.$$
\end{restatable}
In the recent literature (see e.g. \citet{zhang2020adaptive,Puchkin2023,sadiev2023high}), a commonly made assumption is that $\E{\|\vec{g} - \nabla \ell (\vec{w})\|^p} \leq \sigma^2$ for some $p < 2$ where $\sigma^2 >0$ cannot depend on $\vec{w}$. On the contrary, \cref{Ass:oracle with linear stable noise} can accommodate much stronger noise. To illustrate this, consider $d=1$ and $\ell(\vec{w}) = \vec{w}^2/2$. \cref{Ass:oracle with linear stable noise} allows the noise model $\vec{\Sigma}(\vec{w})\vec{\zeta}$ with $\vec{\Sigma}(\vec{w}) = \sqrt{c_1+c_2\vec{w}^2}$, which results in strong noise, since $\vec{\zeta}$ has infinite variance if $\alpha<2$, 
and this is amplified by $\vec{\Sigma}(\vec{w})$ when $\vec{w}$ is large. Such strong noise makes vanilla \SGD{} diverge, as we will illustrate in our experiments.


\begin{restatable}{corollary}{corstableshort}\label{T:CV SMGD stable noise}
Let $n \geq 3$ be odd and $p=1$. 
Consider the iterates from \eqref{eqn:approx-median-descent} with $\vec{e}_t = 0$, and suppose that $\bg_t^{(i)} = \nabla \ell(\vec{w}_t) + \vec{\Sigma}(\bw_t)\vec{\zeta}_t^{(i)}$ verifies \cref{Ass:oracle with linear stable noise}. 
% 
\begin{enumerate}[topsep=0pt,label=(\roman*)]
    \item If $\ell$ is $L$-smooth and $\eta = \mathcal{O}(\tfrac{1}{\sqrt{T}})$, then we have that $  \min_{0\leq t\leq T-1} \E{\|\nabla \ell(\bw_t)\|^2} = \mathcal{O}(1/\sqrt{T}) $. 
    \item If $\ell$ is convex and Lipschitz and $\eta = 1/\sqrt{T}$ then we have that $\mathbb{E}\left[ \ell(\bar \bw_T) - \inf \ell \right] =  \mathcal{O}(1/\sqrt{T}) $, where $\bar \bw_T$ is defined as in \cref{P:CV SMGD convex lipschitz}.
\end{enumerate}
\end{restatable}

This result illustrates how powerful the sample median can be when used in gradient-based optimization.
% 
However, \SMGD{} \eqref{eqn:approx-median-descent} has clear drawbacks: (i) it requires to draw multiple samples per iteration (multiple calls of the oracle), and (ii) we need to compute the sample median up to some tolerance. For instance, the geometric median can not be computed in closed form and there exists a vast body of literature on how to approximate it~\citep{Cohen2016}. Moreover, the noise still needs to be symmetric in order for the median to denoise correctly (cf.\ \cref{Ass:oracle with linear stable noise} and \cref{fig:toy-example-ablation} (right)).

Nevertheless, our convergence guarantees for \SMGD{} allow us to be optimistic that the strong performance in the heavy-tailed setting will also transfer to the online drop-in replacements (of \eqref{eqn:approx-median-descent}), which we had previously presented in \cref{sec:grad-est-in-the-wild}. We investigate this in the following  experiments section.

%
\begin{figure}
    \centering
    \begin{subfigure}[b]{0.49\columnwidth}  
        \includegraphics[width=0.99\textwidth]{{figures/toy_example/moments_alpha1.1}.pdf}
        \caption{$\alpha=1.1$}
    \end{subfigure}
    \begin{subfigure}[b]{0.49\columnwidth}  
        \includegraphics[width=0.99\textwidth]{{figures/toy_example/moments_alpha1.5}.pdf}
        \caption{$\alpha=1.5$}
    \end{subfigure}
    \captionof{figure}{First and second moment of sample median/mean for standard $\alpha$-stable distribution. We approximate $\mathbb{E}$ by running $10\,000$ trials. The sample median has much smaller variance (cf.\ \eqref{eqn:sample-median-noise-cond}) compared to the sample mean.}
    \label{fig:moment-sample-median-mean}
\end{figure}



%%%%%%%%%%%%%%%%%%%%%%%%%%%%%%%%%%%%%%%%%%%%%%

\section{Experiments}\label{sec:experiments}
%
We compare the iterative gradient estimation techniques for two experimental setups, namely (i) estimating a vector/gradient from a fixed distribution, and (ii) methods of the form \eqref{eq:online-update}, that is, iterative gradient estimation while simultaneously learning the weights.

We study both settings for synthetic data where we can control the level of heavy-tailedness, as well as for language modelling tasks with transformer models.  We consider the same three language modeling tasks as studied in \citep{kunstner2023noise}: an encoder-only transformer for the \texttt{PTB} dataset, a Transformer-XL model for the \texttt{WikiText-2} dataset, and fine-tuning a \texttt{DistillBERT} model for question-answering on the \texttt{SQuAD} dataset.
We refer to \cref{sec:app-experiments-supp} for details.

In the following, we use \VC{} to refer to vectorwise clipping~\eqref{eq:vclip-update} and \CC{} for componentwise clipping~\eqref{eq:cclip-update}.
%
\subsection{Gradient Estimation with Fixed Weights}\label{sec:exp-fixed-weight}
%
\paragraph{Synthetic data.}
In the first experiment, we verify the hypothesis that the choice of $\dist$ matters for problem \eqref{prob:gradient-learning-3}; in particular, we show that estimating the median is more stable when the noise distribution is heavy-tailed. 

\emph{Setup.} We consider the problem of estimating a fixed vector $\vec{\hat{g}}$
from the oracle $\vec{g}\sim \vec{\hat{g}} +\vec{\xi}$ with $\vec{\xi}\sim \mathcal{P}$, under varying degree of heavy-tailedness of $\mathcal{P}$. For this purpose, we choose the standard, unskewed $\alpha$-stable distribution $\mathcal{P}_{\alpha}$.
We generate $\vec{\hat{g}}\in \R^d$ with $d=10$ where each component is generated i.i.d standard Gaussian. 
In each iteration, a sample $\bg_t$ is generated as follows: each coordinate of $(\bg_t)_i$ is (independently) sampled from $\mathcal{P}_{\alpha}$ with location $(\hat \bg)_i$ and varying values for $0<\alpha\leq 2$.
Importantly, the median of $\mathcal{P}_{\alpha}$ is equal to the location (i.e.\ $\vec{\hat{g}}$) for all values of $\alpha$. Recall that for $\alpha \leq 1$ the mean of $\mathcal{P}_{\alpha}$ is not defined, and otherwise equal to the median.

We run the \SPP{} method \eqref{eqn:online-spp-update} for $\dist\in \{\tfrac12 \|\cdot\|^2,\|\cdot\|_2,\|\cdot\|_1 , H_{\mu} \}$, with a fixed step size $\tau =0.01$ and track the relative $\ell_2$-error $\frac{\|\vec{m}_t - \vec{\hat{g}}\|_2}{\|\vec{\hat{g}}\|_2}$. 
The corresponding methods are called momentum, \VC{}, \CC{}, and \HC{}, where we set $\mu=1.345$ according to \citet{Huber1981}.
For each method and distribution, we run $1000$ iterations and $50$ different seeds.

\emph{Discussion.} From \cref{fig:toy-example} we find that as the noise becomes more heavy-tailed (as $\alpha$ decreases), momentum fails to produce accurate estimates of $\vec{\hat{g}}$. On the other hand, both \VC{} and \CC{} are robust to heavy tails, as expected, as we showed that \VC{} and \CC{} are estimating the median. Notably \CC{} becomes more accurate (relatively) as $\alpha$ decreases. The Huber function produces results very similar to \VC{}, but slightly inferior. We refer to \cref{fig:toy-example-ablation} for additional ablations on the choice of $\tau$ and skewed noise.

%
\begin{figure}
    \centering
    \begin{subfigure}[b]{0.38\columnwidth}  
        \includegraphics[width=0.99\textwidth]{figures/toy_example/levy_all_noise.pdf}
        %\caption{}
    \end{subfigure}
    \begin{subfigure}[b]{0.55\columnwidth}  
        \includegraphics[width=0.99\textwidth]{figures/toy_example/levy_convergence.pdf}
        %\caption{}
    \end{subfigure}
    \caption{$\tau=0.01$ \textbf{Left:} Final error for varying values of $\alpha$ (from left to right, distributions are more heavy-tailed). Shaded area marks minimal and maximal value over the $50$ independent runs. 
    \textbf{Right:} Convergence plot for all methods for $\tfrac{1}{\alpha}\in[0.5,1.5]$ (higher value of $\tfrac{1}{\alpha}$ corresponds to heavier tails).}
    \label{fig:toy-example}
\end{figure}


\paragraph{Language modelling.}
We also run momentum, \VC{} and \CC{} in order to estimate the full gradient for the language modeling tasks, where the weights are \emph{fixed at initialization}. 
The setup is identical to the one described by \citet{kunstner2023noise}, in particular Figure 1 therein. 
We present the results in \cref{sec:app-fixed-weight-transformer}.
 

%%%%%%%%%%%%%%%%%%%%%%%%%%%%%%%%%%%%%%%%%%%%%%%%%%%%%%
%%%%%%%%%%%%%%%%%%%%%%%%%%%%%%%%%%%%%%%%%%%%%%%%%%%%%%

\subsection{Training with Online Gradient Estimation}\label{sec:experiments-training}

We now present comparisons of various (online) gradient estimators when simultaneously training the weights, that is the distribution of gradients changes in each iteration.

\paragraph{Least squares with heavy-tailed noise.} We consider the noise-perturbed least-squares problem $\min_{\vec{w}\in \R^d} \tfrac12 \| \vec{w} \|^2$
% \begin{equation}
%     ,
% \end{equation}
and use a gradient oracle given by $\vec{g} \sim \vec{w} + \vec{\xi} \in \R^d$, where $\vec{\xi} \sim \mathcal{P}$ is the noise term.
We consider three setups:
\begin{enumerate}[label=(S\arabic*)]
    \item \label{item:independent} The $d$ components of $\vec{\xi}$ are i.i.d.\ distributed according to a standard $\alpha$-stable distribution with location zero and $\alpha=1.1$. Note that $\EE{\vec{\xi} \sim \mathcal{P}}{\vec{\xi}} = \vec{0}$.
    \item \label{item:independent-state} Same as \ref{item:independent}, but the noise is state-dependent, given by $\vec{\xi}_t = \sigma \cdot \vec{\xi}$ where $\sigma = \sqrt{1+\|\vec{w}_t\|^2}$.
    \item \label{item:dependent} $\vec{\xi}$ follows an elliptically contoured $\alpha$-stable distribution \citep{Nolan2013} with location zero and $\alpha=1.1$. In particular, the components of $\vec{\xi}$ are not independent.
\end{enumerate}

We run seven different methods: on the one hand, the method depicted in \cref{sec:sample-median-methods}, \eqref{eqn:approx-median-descent}, where in each iteration we draw $n=5$ samples. We run three variations of \eqref{eqn:approx-median-descent}, namely using the $\ell_1$- and $\ell_2$-sample median as well as the sample mean for $\vec{m}_t$. The $\ell_2$-median is approximately computed using the method proposed in \citet{Vardi2000}. 

\begin{figure}
    \centering
    \includegraphics[width=0.7\columnwidth]{figures/least_squares/benchmark/legend.pdf}
    \begin{subfigure}[b]{0.32\columnwidth}  
        \includegraphics[width=0.99\textwidth]{{figures/least_squares/benchmark/alpha_1.1_independent_True_state_dependent_False_scale_1.0}.pdf}
        \caption{Setting \ref{item:independent}}
    \end{subfigure}
    \begin{subfigure}[b]{0.32\columnwidth}  
        \includegraphics[width=0.99\textwidth]{{figures/least_squares/benchmark/alpha_1.1_independent_True_state_dependent_True_scale_1.0}.pdf}
        \caption{Setting \ref{item:independent-state}}
    \end{subfigure}
    \begin{subfigure}[b]{0.32\columnwidth}  
        \includegraphics[width=0.99\textwidth]{{figures/least_squares/benchmark/alpha_1.1_independent_False_state_dependent_False_scale_1.0}.pdf}
        \caption{Setting \ref{item:dependent}}
    \end{subfigure}
    \caption{Least squares with heavy-tailed noise that is independent over components (\textbf{left}), state-dependent (\textbf{middle}), and where components of noise are dependent (\textbf{right}). 
    All methods in gray use $n=5$ samples per iteration, all others only one.
    Note that for \ref{item:independent-state}, the sample mean \textbf{immediately diverges}, and hence does not appear in the plot.
    Shaded area depicts minimum and maximum value over 50 repetitions.}
    \label{fig:least-squares-example}
\end{figure}

On the other hand, we run our online median estimator methods \eqref{eq:online-update}, with $\dist \in \{\|\cdot\|_2,~ \|\cdot\|_1, ~H_\mu\}$, that is, \VC{}, \CC{}, and \HC{}. Importantly, these methods only receive \emph{one} sample per iteration, and they do not involve any median computations (in contrast to $\ell_1$-  and $\ell_2$-sample median).
We also run \texttt{clipped-SGD} (cf.\ \eqref{eqn:clipped-sgdm}) with $\beta=0.9$ and $c=50$.\footnote{This choice is a tradeoff between slow convergence (large $c$) and high instability (small $c$).} We run all methods with a learning rate $\eta_t=0.01$. For \VC{}, \CC{}, and \HC{} we set $\tau=1$, and we set again $\mu=1.345$ for \HC{}. We always run 50 repetitions and report averaged metrics.

\emph{Discussion.} \cref{fig:least-squares-example} shows that both $\ell_1$- and $\ell_2$-sample median attain the smallest objective, while the sample mean does not converge as predicted by our theory (see discussion after \cref{prop:approx-median-descent}). From the online methods, \CC{} performs best in the scenario where the noise is componentwise independent, while \VC{} and \HC{} are better in the other case. Overall, \HC{} and \VC{} behave again very similarly (a more detailed analysis of \HC{} with different values of $\mu$ can be found in \cref{fig:huber-mu-ablation}). We find that both \VC{} and \CC{} reach much smaller loss values than clipping the stochastic gradient directly, and thus clipping only the increment is better suited here. We also find that the $\ell_1$-sample median can deal with state-dependent noise \ref{item:independent-state}  as predicted by our theory. 
%

\paragraph{Language modeling.} We now consider the online methods \eqref{eq:online-update}, that is training while simultaneously estimating the gradient, for the three language modeling tasks.
We compare \SGDM{}, \VC{} and \CC{}, and \texttt{clipped-SGD}.\footnote{We also tried \HC{} (with $\mu=1.3$) on these problems, but did not observe an advantage over \VC{} (plots not shown).} We also add \texttt{Adam} as a baseline, but focus on the comparison among the \SGD{}-type methods. For all methods, we tune the learning rate on a $\log_{10}$-scaled grid (tuned values reported in \cref{tab:transformer-lr-values}), displaying the one that attained smallest final training loss (averaged over $5$ seeds). 

We choose a standard momentum/clipping parameters for all tasks (without tuning): we set~$\beta=0.9$ for \SGDM{},~$\tau=0.1$ for \texttt{V/CClip}, and~$\beta=0.9,~c=1$ for \texttt{clipped-SGD}.

\emph{Discussion.} \cref{fig:transformer-training} shows that \VC{} improves over \SGDM{} for \texttt{PTB} and \texttt{WikiText-2}, and is on par for \texttt{SQuAD}. \CC{} performs always worse, which might be due to the noise being not componentwise independent.
However, \VC{} does not close the gap to \texttt{Adam}, and also performs worse than \texttt{clipped-SGD} for the \texttt{WikiText-2} dataset.


\begin{figure}
    \centering
    \includegraphics[width=0.7\columnwidth]{figures/transformer_example/trainer/legend.pdf}
    \begin{subfigure}[b]{0.33\columnwidth}  
        \includegraphics[width=0.99\textwidth]{{figures/transformer_example/trainer/ptb_transformer_encoder_L/train_loss}.pdf}
        \caption{\texttt{PTB}}
    \end{subfigure}
    \begin{subfigure}[b]{0.3\columnwidth}  
        \includegraphics[width=0.99\textwidth]{{figures/transformer_example/trainer/wikitext2_transformer_L/train_loss}.pdf}
        \caption{\texttt{WikiText-2}}
    \end{subfigure} 
    \begin{subfigure}[b]{0.3\columnwidth}  
        \includegraphics[width=0.99\textwidth]{{figures/transformer_example/trainer/distillbert_squad_L/train_loss}.pdf}
        \caption{\texttt{SQuAD}}
    \end{subfigure}
    \caption{Training loss for each method, with tuned learning rate. Shaded area depicts minimum and maximum value over 5 seeds.}
    \label{fig:transformer-training}
\end{figure}
%
%
%
\begin{figure}
    \centering
    \includegraphics[width=0.7\columnwidth]{figures/transformer_example/trainer/legend.pdf}
    \begin{subfigure}[b]{0.3\columnwidth}  
        \includegraphics[width=0.99\textwidth]{{figures/transformer_example/trainer/ptb_transformer_encoder_L/valid_ppl}.pdf}
        \caption{\texttt{PTB}}
    \end{subfigure}
    \begin{subfigure}[b]{0.3\columnwidth}  
        \includegraphics[width=0.99\textwidth]{{figures/transformer_example/trainer/wikitext2_transformer_L/valid_ppl}.pdf}
        \caption{\texttt{WikiText-2}}
    \end{subfigure} 
    \begin{subfigure}[b]{0.3\columnwidth}  
        \includegraphics[width=0.99\textwidth]{{figures/transformer_example/trainer/distillbert_squad_L/valid_exact_f1}.pdf}
        \caption{\texttt{SQuAD}}
    \end{subfigure}
    \caption{Validation score (measured as perplexity, or F1 score) for each method, with tuned learning rate. Shaded area depicts minimum and maximum value over 5 seeds.}
    \label{fig:transformer-training-valid}
\end{figure}
%%%%%%%%%%%%%%%%%%%%%%%%%%%%%%%%%%%%%%%%%%


\section{Conclusion}

We provide a new approach for robust gradient estimation using the stochastic proximal point method, with a focus on (online) estimators of the median gradient.
We observe that this general framework intersects with many existing works on robust stochastic optimization, where the connection between clipping and median estimation was not known previously. 
%

We provide an analysis of \SGD{} using the median of sampled gradients (called \SMGD{}), and show convergence guarantees under potentially heavy-tailed and state-dependent noise. In contrast, under the same assumption, \SGD{} with the sample mean does not converge.

Finally, numerical experiments show that online versions of \SMGD{}, even though not covered by this theory, perform equally well on problems with heavy-tailed noise, while methods based on mean estimation fail. However, when training transformer models on language tasks, this gap is far less pronounced; this could indicate that the level of heavy-tailedness in these applications is moderate.

\subsubsection*{Acknowledgments}
We thank the Scientific Computing Core at the Flatiron Institute, a division of the Simons Foundation, for the compute facilities and support.
% 
U.\c{S}. and F.S. are partially supported by the French government under the
management of Agence Nationale de la Recherche as part of the “Investissements d’avenir” program,
reference ANR-19-P3IA-0001 (PRAIRIE 3IA Institute), and the
European Research Council Starting Grant DYNASTY – 101039676.
G.G. gratefully acknowledges the Flatiron Institute for their hospitality and financial support during the visit when this work was written.


\clearpage


\bibliography{references}
\bibliographystyle{tmlr}

\clearpage

\tableofcontents

\appendix

\section*{Appendix}

%
\section{Missing Proofs for \cref{sec:sample-median-methods}}\label{sec:missing-proofs}

\subsection{On the Assumptions on the Sample Median}

We first state a simple fact about Lipschitz functions:
\begin{lemma}\label{L:oracle purely bouded if lipschitz objective}
    If \cref{ass:uniform-bias-variance-bound-median} holds, and if $\ell$ is Lipschitz continuous, then without loss of generality we can assume that $\delta_2 = \sigma_2=0$.
\end{lemma}

\begin{proof}
    If $\ell$ is $G$-Lipschitz, then for every subgradient $\vec{u} \in \partial \ell(\bw)$ we have $\Vert \vec{u} \Vert \leq G$, so one can replace $\delta_1^2$ with $\delta_1^2 + \delta_2 G^2$ and $\sigma_1^2$ with $\sigma_1^2 + \sigma_2 G^2$.
\end{proof}


In the following proof of \cref{prop:approx-median-descent} we will need an equivalent but more compact formulation of \cref{ass:uniform-bias-variance-bound-median}.
%
\begin{lemma}\label{lem:noise-assum-equiv}
    Let $\m$ be a random variable in $\mathbb{R}^d$, and consider the following property:
    \begin{equation}\label{ass:approx-median-bounded-variance}
        \E{\Vert \vec{m} - \nabla \ell(\vec{w}) \Vert^2} \leq \nu_1^2 + \nu_2 \|\nabla \ell(\bw)\|^2.
    \end{equation}
     If \eqref{ass:noise-approx-median} holds then \eqref{ass:approx-median-bounded-variance} holds with $\nu_1^2 = \delta_1^2+\sigma_1^2 $ and $\nu_2=\delta_2+\sigma_2$.
     If \eqref{ass:approx-median-bounded-variance} holds, then \eqref{ass:noise-approx-median} holds with $\delta_1=\sigma_1 = \nu_1$ and $\delta_2=\sigma_2=\nu_2$.
\end{lemma}
%
\begin{proof}
    Expanding the squares, we get
    \begin{equation*}
        \E{\Vert \vec{m} - \nabla \ell(\vec{w}) \Vert^2}
        =
        \E{\Vert \vec{m} - \E{\vec{m}} \Vert^2}
        +
        \Vert  \E{\vec{m}} - \nabla \ell(\vec{w}) \Vert^2,
    \end{equation*}
    where we used that $\E{\langle \vec{m} - \E{\vec{m}}, \E{\vec{m}} - \nabla \ell(\vec{w}) \rangle} = 0$.
    Now, the first implication is trivial, and the second follows from the positivity of the involved terms.
\end{proof}

We now verify that \cref{ass:uniform-bias-variance-bound-median} holds true for the $\ell_1$-sample median, provided that the noise on the subgradient is linearly depending on an $\alpha$-stable distribution, and that the linear coefficients grow with the norm of the subgradient.



\begin{restatable}{lemma}{propnoisecond}\label{stable noise implies finite bias variance v2}
Let $n \in \mathbb{N}$ be odd and $n \geq 3$. Suppose that the subgradient oracles $\bg^{(i)} = \nabla \ell(\vec{w}) + \vec{\Sigma}(\bw)\vec{\zeta}^{(i)}$ verify Assumption \ref{Ass:oracle with linear stable noise}, for all $i=1,\dots,n$ and set $$\vec{m} = \samplemedian{\ell_1}\left(\bg^{(i)}(\vec{w})\right).$$
% 
Then, \cref{ass:uniform-bias-variance-bound-median} holds with $\delta_1 = \delta_2 =0$, and  $\sigma_1^2 = C_n c_1$, $\sigma_2 = C_n c_2$, for some constant $C_n = \mathcal{O}(1/n)$.
\end{restatable}

\begin{proof}
    
Let us denote the noise vector by $\vec{\xi}^{(i)} = \vec{\Sigma}(\bw)\vec{\zeta}^{(i)}$. By the stability property of $\alpha$-stable distributions \citep[Proposition 1.3]{Nolan2020}, we have that \begin{align}
    \label{eqn:marg_dist}
   {\xi}_k^{(i)}: =\vec{\xi}_k^{(i)}(\vec{w}) = \sum_{j=1}^d [\vec{\Sigma}(\vec{w})]_{kj} \vec{\zeta}^{(i)}_j =^{\mathrm{d}}  \|\vec{\Sigma}(\vec{w})_k\|_{\alpha} z ,
\end{align}
where $=^{\mathrm{d}}$ denotes equality in distribution, $z$ is a standard symmetric $\alpha$-stable random variable in $\mathbb{R}$, and $\vec{\Sigma}(\vec{w})_k$ denotes the $k$-th row of $\vec{\Sigma}(\vec{w})$.
Observe that from \cref{Ass:oracle with linear stable noise} we can write that
\begin{equation}\label{eqn:frobenius to alpha norm}
    \sum_{i=1}^d \|\vec{\Sigma}(\vec{w})_i\|_{\alpha}^2 
    \leq
    \rho_\alpha^2\sum_{i=1}^d \|\vec{\Sigma}(\vec{w})_i\|_{2}^2 
    =
    \rho_\alpha^2\Vert \vec{\Sigma}(\vec{w}) \Vert_F^2
    \leq 
    c_1' + c_2' \|\nabla \ell(\vec{w})\|^2, \quad \mbox{for all  } \vec{w} \in \R^d.
\end{equation}
by setting $c_1' =\rho_\alpha^2 c_1$ and $c_2 = \rho_\alpha^2 c_2$, and where $\rho_\alpha>0$ is the best constant such that $\Vert \cdot \Vert_\alpha \leq \rho_\alpha \Vert \cdot \Vert_2$, which is $\rho_\alpha = d^{\tfrac{2-\alpha}{2\alpha}}$ is for $\alpha \leq 2$, and $\rho_\alpha=1$ otherwise.

We start by computing the bias term that is required by \cref{ass:uniform-bias-variance-bound-median}.
\begin{align*}
   \Vert \E{\vec{m}} - \nabla \ell(\vec{w}) \Vert^2 =& \Vert \E{ \samplemedian{\ell_1}\left(\nabla \ell(\vec{w}) + \vec{\xi}^{(i)}(\vec{w})\right) } - \nabla \ell(\vec{w}) \Vert^2 \\
   =& \Vert \E{ \samplemedian{\ell_1}\left(\vec{\xi}^{(i)}(\vec{w}) \right) }  \Vert^2 % \\
   % =& 0,
\end{align*}
where we used the fact that 
$$\samplemedian{\ell_1}\left( \nabla \ell(\vec{w}) + \vec{\xi}^{(i)}(\vec{w}) \right) = \nabla \ell(\vec{w}) + 
\samplemedian{\ell_1}\left( \vec{\xi}^{(i)}(\vec{w})\right).$$
% 
On the other hand, the $\ell_1$-median boils down to computing component-wise medians, hence we have that
\begin{align*}
    \E{\samplemedian{\ell_1}\left( \vec{\xi}^{(i)}(\vec{w})\right) }= \begin{bmatrix}
\E{\median_{i\in[n]}\left( \xi_1^{(i)}\right)  } \\
\vdots \\
\E{\median_{i\in[n]}\left( \xi_d^{(i)}\right) }
\end{bmatrix},
\end{align*}
where $\median_{i\in[n]}(\xi_k^{(i)}) = \argmin{m \in \mathbb{R}} \frac1{n} \sum_{i=1}^n |m-\xi_k^{(i)}| $ is the one dimensional sample median estimator. 

Now, let us compute $\E{\median_{i\in[n]}\left( \xi^{(i)}_k\right) }$ for some $k \in \{1,\dots, d\}$. Denoting $r = (n-1)/2 \in\mathbb{N}$ and by using the probability density function of the sample median (cf.\ \citet{maritz1978note}), we have that:
\begin{align*}
    \E{\median_{i\in[n]}\left( \xi^{(i)}_k\right) } =& \frac{n!}{r!r!} \underbrace{\int_{\mathbb{R}} x \left( F_{\vec{w},k}(x) (1-  F_{\vec{w},k}(x)) \right)^r p_{\vec{w},k}(x) \mathrm{d} x}_{=:\mathrm{A}},
\end{align*}
where $F_{\vec{w},k}$ and $p_{\vec{w},k}(x)$ are respectively the cumulative distribution function (cdf) and the probability density function (pdf) of $\xi_k^{(i)} \sim \mathcal{P}_{\alpha, \|\vec{\Sigma}_{k}(\vec{w})\|_{\alpha} }$ that is a symmetric $\alpha$-stable random variable with scale $\|\vec{\Sigma}_{k}(\vec{w})\|_{\alpha} $. As $p_{\vec{w},k}$ is continuous and symmetric around zero, we have $p_{\vec{w},k}(x) = p_{\vec{w},k}(-x)$ and $ (1-  F_{\vec{w},k}(x)) = F_{\vec{w},k}(-x)$ for all $x\in \mathbb{R}$. This gives us:
\begin{align*}
    % \E{\median_{i\in[n]}\left( \xi^{(i)}_k\right) }  
    \mathrm{A} =& \int_{0}^\infty x \left( F_{\vec{w},k}(x)  F_{\vec{w},k}(-x) \right)^r p_{\vec{w},k}(x) \mathrm{d} x + \int_{-\infty}^0 x \left( F_{\vec{w},k}(x)  F_{\vec{w},k}(-x) \right)^r p_{\vec{w},k}(x) \mathrm{d} x \\
    =& \int_{0}^\infty x \left( F_{\vec{w},k}(x)  F_{\vec{w},k}(-x) \right)^r p_{\vec{w},k}(x) \mathrm{d} x - \int_{0}^\infty x \left( F_{\vec{w},k}(x)  F_{\vec{w},k}(-x) \right)^r p_{\vec{w},k}(x) \mathrm{d} x \\
    =& 0.
\end{align*}
Hence, we obtain that
$$\E{ \samplemedian{\ell_1}\left( \vec{\xi}^{(i)}(\vec{w}) \right) } = 0.$$
This shows that the first line of \cref{ass:uniform-bias-variance-bound-median} holds with $\delta_1 = \delta_2 =0$. 

We now proceed to the estimation of the variance. By using similar arguments, we have that:
\begin{align*}
   \E{ \Vert \vec{m} - \E{\vec{m} } \Vert^2  } =\; & \E{ \Vert  \samplemedian{\ell_1}\left(\vec{\xi}^{(i)}(\vec{w}) \right)  \Vert^2  } \\
   = \;& \sum_{k=1}^d \E{ \left(  \median_{i \in [n]}\left( \xi_k^{(i)} \right)  \right)^2  } \\
   \numberthis \label{eqn:var_median}
   =\; & \sum_{k=1}^d \frac{n!}{r!r!}  \mathbb{E}\left[\Bigl(\xi_k^{(i)}\Bigr)^2 \Biggl(F_{\vec{w},k}\Bigl(\xi_k^{(i)}\Bigr)\Bigl(1-F_{\vec{w},k}\Bigl(\xi_k^{(i)}\Bigr)\Bigr) \Biggr)^r \right],
\end{align*}
where \eqref{eqn:var_median} follows from \citet[Equation 2.1]{maritz1978note}. By using \eqref{eqn:marg_dist} with \eqref{eqn:frobenius to alpha norm} and denoting the cdf of the standard symmetric $\alpha$-stable distribution by $F_{z}$, we further compute that:
\begin{align*}
   \E{ \Vert \vec{m} - \E{\vec{m}} \Vert^2 } 
   = \; & \frac{n!}{r!r!} \mathbb{E}\left[z^2  \Biggl(F_{z}(z)(1-F_{z}(z) \Biggr)^r  \right]   \sum_{k=1}^d \|\vec{\Sigma}_{k}(\vec{w})\|_{\alpha}^2  \\
   \leq \; & \frac{n!}{r!r!} \mathbb{E}\left[z^2  \Biggl(F_{z}(z)(1-F_{z}(z) \Biggr)^r  \right] \left(c_1' + c_2' \|\nabla \ell(\vec{w})\|^2 \right) \\
   =: \; & C_n \left(c_1' + c_2' \|\nabla \ell(\vec{w})\|^2 \right).
\end{align*}
This shows that the second line of \cref{ass:uniform-bias-variance-bound-median} holds with $\sigma_1^2 = C_n \rho_\alpha^2 c_1$ and $\sigma_2 = C_n \rho_\alpha^2 c_2$.


Finally we observe that, again by \citet[Equation 2.1]{maritz1978note}, $C_n$ is the variance of the sample median of a set of i.i.d. standard symmetric $\alpha$-stable variables $\{z_1, \dots, z_{n}\}$, where $z_i =^{\mathrm{d}} z$. As $n\to \infty$, it is well-known that the sample median is asymptotically normal with variance $1/(4n f^2(0))$, where $f$ denotes the pdf of $z$ (see \citet{maritz1978note}). Hence we conclude that $C_n = \mathcal{O}(1/n)$ as $n\to \infty$. This concludes the proof.
\end{proof}



\subsection{Rates for \SMGD{} for Smooth Nonconvex Functions: Proof of \cref{prop:approx-median-descent} }

The following proposition is an application of biased \SGD{} results to the \SMGD{} algorithm with errors.
It is possible to show that \cref{ass:uniform-bias-variance-bound-median,Ass:errors bounded expectation} both imply Assumption 9 from \citet{DemMalSokRic23}.
So from a qualitative point of view the result below is the same that what we would obtain by using Theorem 3 from \citet{DemMalSokRic23}.
But doing so would produce a bound with worse constants, so we prefer to provide a proof exploiting directly our assumptions.


%%%%%%%%%%%%%%%%%%%%%
\propapproxmed*
%%%%%%%%%%%%%%%%%%%%%
\begin{proof}
    In this proof we note $\mathbb{E}_t$ for the expectation taken conditionally to the filtration generated by $\bw_0, \dots \bw_t$.
    In particular, we will have from \cref{Ass:errors bounded expectation} on the errors together with the tower rule  that $\mathbb{E}_{t} \Vert \be_t \Vert^2 = \mathbb{E}_{t} \mathbb{E}\left[ \Vert \be_t \Vert^2 \ | \ \bg_t^{(1)}, \dots, \bg_t^{(n)} \right]  \leq \varepsilon^2$. 

    
    Smoothness of $\ell$ implies
    \begin{equation*}
        \ell(\bw_{t+1}) - \ell(\bw_t) - \langle \nabla \ell(\bw_t), \bw_{t+1} - \bw_t \rangle \leq \frac{L}{2}\Vert \bw_{t+1} - \bw_t \Vert^2.
    \end{equation*}
    If we note $\vec{d}_t := \vec{m}_t + \vec{e}_t$, thus $\bw_{t+1} - \bw_t = -\eta \vec{d}_t$. This implies
    \begin{equation*}
        \ell(\bw_{t+1}) - \ell(\bw_t) + \eta \langle \nabla \ell(\bw_t), \vec{d}_t \rangle \leq \frac{L\eta^2}{2}\Vert \vec{d}_t \Vert^2.
    \end{equation*}
    Let us bound the two terms depending on $\vec{d}_t$.
    First, take expectation conditioned to $\mathcal{F}_t$ (which we will denote $\mathbb{E}_t$) to write
    \begin{eqnarray*}
        \EE{t}{\langle \nabla \ell(\bw_t), \vec{d}_t \rangle}
        & = &
        \langle \nabla \ell(\bw_t), \EE{t}{\vec{d}_t} \rangle
        =
        \frac{1}{2}\Vert \nabla \ell(\bw_t) \Vert^2 + \frac{1}{2}\Vert \EE{t}{\vec{d}_t} \Vert^2 - \frac{1}{2}\Vert \EE{t}{\vec{d}_t} - \nabla \ell(\bw_t) \Vert^2 \\
        & \geq &
        \frac{1}{2}\Vert \nabla \ell(\bw_t) \Vert^2 - \frac{1}{2}\Vert \EE{t}{\vec{d}_t} - \nabla \ell(\bw_t) \Vert^2.
    \end{eqnarray*}
    Using \eqref{ass:noise-approx-median} and $\Vert \EE{t}{\vec{e}_t} \Vert^2\leq \EE{t}{\|\vec{e}_t\|^2}\leq \varepsilon^2$ due to Jensen's inequality, we can bound
    \begin{equation*}
        \Vert \EE{t}{\vec{d}_t} - \nabla \ell(\bw_t) \Vert^2
        \leq 
        2 \Vert \EE{t}{\vec{m}_t} - \nabla \ell(\bw_t) \Vert^2
        +
        2 \Vert \EE{t}{\vec{e}_t} \Vert^2
        \leq 2\delta_1^2 + 2\delta_2 \|\nabla \ell(\bw_t)\|^2+ 2 \varepsilon^2.
    \end{equation*}
    Second, we use Young's inequality and our assumptions to bound $\Vert \vec{d}_t \Vert^2$:
    \begin{eqnarray*}
        \EE{t}{\Vert \vec{d}_t \Vert^2}
        & \leq &
        2\EE{t}{\Vert \vec{e}_t + \vec{m}_t - \nabla \ell(\bw_t) \Vert^2}
        +2 \Vert \nabla \ell(\bw_t) \Vert^2 \\
        &\leq & 
        4 \EE{t}{ \Vert \vec{e}_t \Vert^2} + 4 \EE{t}{\Vert \vec{m}_t - \nabla \ell(\bw_t) \Vert^2} +2 \Vert \nabla \ell(\bw_t) \Vert^2 \\
        & \leq & 
        4 \varepsilon^2 + 4(\delta_1^2+\sigma_1^2 ) + (2+4\delta_2+4\sigma_2) \Vert \nabla \ell(\bw_t) \Vert^2,
    \end{eqnarray*}
    where in the last inequality we used that \eqref{ass:noise-approx-median} implies \eqref{ass:approx-median-bounded-variance} (cf. \cref{lem:noise-assum-equiv}).
    Combine all the above inequalities to get
    \begin{align*}
        &\EE{t}{\ell(\bw_{t+1})} - \ell(\bw_t)
        + \tfrac{\eta}{2}\Vert \nabla \ell(\bw_t) \Vert^2
        - \tfrac{\eta}{2}(2\delta_1^2 + 2\delta_2 \|\nabla \ell(\bw_t)\|^2 +\varepsilon^2) \\  
        &\hspace{2ex}\leq \frac{L\eta^2}{2}\big( 
        4(\varepsilon^2 + \delta_1^2 + \sigma_1^2 )+(2+4\delta_2+4\sigma_2)\|\nabla \ell(\bw_t)\|^2\big) .
    \end{align*}
    Rearranging yields
    \begin{align*}
        &\EE{t}{\ell(\bw_{t+1})} - \ell(\bw_t)\\
        &\hspace{3ex}\leq -\eta \big[\tfrac12 - \delta_2 - L\eta(1+2\sigma_2+2\delta_2)\big]
        \|\nabla \ell(\bw_t)\|^2
        + \eta (\delta_1^2+\varepsilon^2) + 2L\eta^2(\varepsilon^2+ \delta_1^2 + \sigma_1^2).
    \end{align*}
    If $\delta_2\leq \tfrac18$ and using $\eta \leq \frac{1}{8L(1+2\sigma_2+2\delta_2)}$, after taking full expectation, we obtain
    \begin{eqnarray*}
        \E{ \ell(\bw_{t+1})} - \E{\ell(\bw_{t})}
        & \leq &
        -\tfrac{\eta}{4} \E{\Vert \nabla \ell(\bw_t) \Vert^2} 
        + 
        \eta(\delta_1^2 + \varepsilon^2) 
        +
        2L\eta^2 (\varepsilon^2+\delta_1^2 + \sigma_1^2 ).
    \end{eqnarray*}
    Sum over $t=0, \dots , T-1$, multiply by $\frac{4}{\eta T}$ and use $\min\limits_{t=0, \dots, T-1}\E{\Vert \nabla \ell(\bw_t) \Vert^2} \leq \tfrac{1}{T}\sum_{t=0}^{T-1}\E{\Vert \nabla \ell(\bw_t) \Vert^2}$ to finally obtain
    \begin{equation*}
        \min\limits_{t=0, \dots, T-1} \E{\Vert \nabla \ell(\bw_t) \Vert^2}
        \leq 
        \frac{4 (\ell(\bw_0) - \inf \ell)}{\eta T}
        +8L\eta(\varepsilon^2+\delta_1^2 + \sigma_1^2 ) + 4(\delta_1^2 + \varepsilon^2).
    \end{equation*}
\end{proof}


\subsection{Rates for Biased \SGD{} for Convex Lipschitz Functions}\label{SS:SMGD:rates for convex lipschitz general}


In this section we prove another result about the complexity of biased \SGD{}, that is \SGD{} with biased estimators of the subgradient.
As far as we know, biased \SGD{} was only studied in the case of strongly convex functions or smooth nonconvex functions \citep{DemMalSokRic23}, or in the particular case of random projections.
In that regard, the next result is a new contribution to the analysis of biased \SGD{}.

Before stating our result, we recall a few standard definitions given that we are working with a nonsmooth convex function. 
Given a convex function $\ell : \mathbb{R}^d \to \mathbb{R}$, we say that $\vec{u}$ is a subgradient of $\ell$ at $\bw$ if
\begin{equation*}
    (\forall \bw' \in \mathbb{R}^d) \quad  \ell(\bw') - \ell(\bw) - \langle \vec{u}, \bw' - \bw \rangle \geq 0.
\end{equation*}
We note $\partial \ell (\bw)$ the set of all subgradients of $\ell$ at $\bw$, that is called the subdifferential of $\ell$ at $\bw$.


\begin{theorem}\label{T:CV SMGD for convex lipschitz}
    Let $\ell : \mathbb{R}^d \to \mathbb{R}$ be a convex $G$-Lipschitz function such that ${\rm{argmin}}~\ell \neq \emptyset$.
    Let $\eta >0$, and note $D:= \Vert \bw_0 - \bw_* \Vert$ for some $\bw_* \in {\rm{argmin}}~\ell$.
    Consider the iterates generated by a biased Stochastic Subgradient Descent method
    \begin{equation*}
        \bw_{t+1} = \bw_t - \eta \bg_t,
    \end{equation*}
    and assume that the estimator $\bg_t$ has a uniformly finite bias and variance (we denote by $\mathcal{F}_t$ the filtration induced by $\bw_0, \dots, \bw_t$):
    \begin{equation*}
        \Vert \EE{}{\bg_t | \mathcal{F}_t} - \vec{u}_t \Vert  \leq \delta,
        \quad
        \EE{}{ \Vert \bg_t - \EE{}{\bg_t | \mathcal{F}_t}\Vert^2 | \mathcal{F}_t}  \leq \sigma^2,
    \end{equation*}
    where $\vec{u}_t$ is any subgradient of $\ell$ at $\bw_t$.
    Then, for every $T \geq 1$
    \begin{equation*}
        \EE{}{\ell(\bar \bw_T) - \inf \ell}
        \leq 
        \frac{2 D^2}{\eta T}
        + \eta \frac{\sigma^2 + (\delta + G)^2}{2} 
        + \eta  \frac{ \delta^2 T}{4}
    \end{equation*}
    where $\bar \bw_T$ is an average of the first $T$ iterates defined by $\bar \bw_T = \frac{1}{\sum_{t=0}^{T-1} \theta^t} \sum_{t=0}^{T-1} \theta^t \bw_t$, $\theta = \frac{T}{T+2}$.
    In particular, if we take a constant step size equal to 
    \begin{equation*}
    \eta = 
        \frac{2\sqrt{2}D}{\sqrt{2(\sigma^2 + (\delta + G)^2)T+\delta^2 T^2 }}
        \quad
        \text{(resp. $\eta = \frac{1}{\delta T + \sqrt{T}}$ )}
    \end{equation*}
    then we can guarantee that 
    \begin{equation*}
        \EE{}{\ell(\bar \bw_T) - \inf \ell}
        \leq 
        %\sqrt{2\eta_1} D \sqrt{\frac{\sigma^2 + (\delta + G)^2}{T} } + \sqrt{2}D \delta.
        2D \sqrt{\frac{\sigma^2 + (\delta + G)^2}{T}+\frac{\delta^2}{2}}
        \quad
        \text{(resp. $\frac{2D^2 + \sigma^2 + \delta^2 + G^2}{\sqrt{T}} + (2D^2 +1)\delta$ ).}
    \end{equation*}
\end{theorem}


\begin{proof}
    Let $w_* \in {\rm{argmin}}~\ell$ and $T \geq 1$ be fixed.
    Develop the squares and use the definition of the algorithm to write
    \begin{eqnarray*}
        \Vert \bw_{t+1} - \bw_* \Vert^2 
        & = & 
        \Vert \bw_t - \bw_* \Vert^2 + 2  \langle \bw_{t+1} - \bw_t, \bw_t -\bw_* \rangle + \Vert \bw_{t+1} - \bw_t \Vert^2 \\
        &=&
        \Vert \bw_t - \bw_* \Vert^2 - 2 \eta \langle \bg_t, \bw_t -\bw_* \rangle + \eta^2 \Vert \bg_t \Vert^2.
    \end{eqnarray*}
    Let $\bar \bg_t$ be the projection of $\EE{}{\bg_t | \mathcal{F}_t}$ onto $\partial \ell(\bw_t)$, the subdifferential of $\ell$ at $\bw_t$, which is a nonempty closed convex set.
    This means that from our assumption on the bias of $\bg_t$ we will have
    \begin{equation*}
        \Vert \mathbb{E}\left[ \bg_t | \mathcal{F}_t \right] - \bar \bg_t \Vert \leq \Vert \mathbb{E}\left[ \bg_t | \mathcal{F}_t \right] - \vec{u}_t \Vert \leq \delta.
    \end{equation*}
    Then, after taking conditional expectation (denoted by $\EE{t}{\cdot}$ instead of $\EE{}{\cdot | \mathcal{F}_t}$) we can write
    \begin{eqnarray*}
        \EE{t}{\Vert \bw_{t+1} - \bw_* \Vert^2}
        & = & 
        \Vert \bw_t - \bw_* \Vert^2
        - 2 \eta \langle \EE{t}{\bg_t}, \bw_t -\bw_* \rangle + \eta^2 \EE{t}{\Vert \bg_t \Vert^2} \\
        & =& 
        \Vert \bw_t - \bw_* \Vert^2
        - 2 \eta \langle \bar \bg_t, \bw_t -\bw_* \rangle 
        + 2 \eta \langle \bar \bg_t -  \EE{t}{\bg_t}, \bw_t -\bw_* \rangle 
        + \eta^2 \EE{t}{\Vert \bg_t \Vert^2} \\
        & \leq &
        \Vert \bw_t - \bw_* \Vert^2
        - 2 \eta (\ell(\bw_t) - \inf \ell)
        + 2 \eta \langle \bar \bg_t -  \EE{t}{\bg_t}, \bw_t -\bw_* \rangle 
        + \eta^2 \EE{t}{\Vert \bg_t \Vert^2},
    \end{eqnarray*}
    where in the last inequality we used the convexity of $\ell$ together with the fact that $\bar \bg_t \in \partial \ell(\bw_t)$.
    The last term can be bounded by using our assumptions on the estimator $\bg_t$, and the fact that $\ell$ has bounded subgradients:
    \begin{equation*}
        \EE{t}{\Vert \bg_t \Vert^2} 
        = \EE{t}{\Vert \bg_t - \EE{t}{\bg_t} \Vert^2} + \Vert \EE{t}{\bg_t} \Vert^2
        \leq 
        \sigma^2 + (\Vert \EE{t}{\bg_t} - \bar \bg_t \Vert + \Vert \bar \bg_t \Vert)^2
        \leq 
        \sigma^2 + (\delta + G)^2.
    \end{equation*}
    Injecting this in the previous inequality we obtain
    \begin{equation*}
        \EE{t}{\Vert \bw_{t+1} - \bw_* \Vert^2} 
        \leq 
        \Vert \bw_t - \bw_* \Vert^2
        - 2 \eta (\ell(\bw_t) - \inf \ell)
        + 2 \eta \langle \bar \bg_t -  \EE{t}{\bg_t}, \bw_t -\bw_* \rangle 
        + \eta^2 \sigma^2 + \eta^2(\delta + G)^2.
    \end{equation*}
Now we introduce a parameter $\varepsilon >0$, and we use Young's inequality together with our assumption on the bias of $\bg_t$ to write
    \begin{equation*}
        2\langle \bar \bg_t -  \EE{t}{\bg_t}, \bw_t - \bw_* \rangle
        \leq 
        \frac{1}{\varepsilon}\Vert \bar \bg_t -  \EE{t}{\bg_t} \Vert^2 + \varepsilon \Vert \bw_t - \bw_* \Vert^2
        \leq 
        \varepsilon^{-1} \delta^2 + \varepsilon \Vert \bw_t - \bw_* \Vert^2.
    \end{equation*}
    Inject this bound in our main inequality
    \begin{equation*}
        \EE{t}{\Vert \bw_{t+1} - \bw_* \Vert^2} 
        \leq 
        (1 + \eta \varepsilon)\Vert \bw_t - \bw_* \Vert^2
        - 2 \eta (\ell(\bw_t) - \inf \ell)
        + \eta^2 \sigma^2 + \eta^2(\delta + G)^2
        + \eta \varepsilon^{-1} \delta^2,
    \end{equation*}
    and after reorganizing the terms, dividing by $2 \eta$ and taking expectation we finally obtain
    \begin{equation*}
        \EE{}{\ell(\bw_t) - \inf \ell}
        \leq 
        \frac{1 + \eta \varepsilon}{2 \eta}\EE{}{\Vert \bw_t - \bw_* \Vert^2}
        - \frac{1}{2\eta}\EE{}{\Vert \bw_{t+1} - \bw_* \Vert^2} 
        + \eta \frac{\sigma^2 + (\delta + G)^2}{2} 
        +  \varepsilon^{-1} \frac{\delta^2}{2}.
    \end{equation*}
    We are now ready to use a weighted telescopic sum argument.
    First, define $\theta := \frac{1}{1 + \eta \varepsilon}$ and $\rho_t := \theta^{t}$.
    Second, multiply our inequality by $\rho_{t+1}$ and then sum for $t=0, \dots, T-1$.
    Observe that the term $(1+ \eta \varepsilon) \rho_{t+1}$ is equal to $\rho_t$ due to our definition, which means that we have a telescoping sum where the terms $\tfrac{\rho_t}{2\eta} \EE{}{\Vert \bw_t - \bw_* \Vert^2}$ will cancel each other.
    Third, divide by $\sum_{t=0}^{T-1} \rho_{t+1}$ so to obtain
    \begin{equation*}
    \frac{1}{\sum_{t=0}^{T-1} \rho_{t+1}}\sum_{t=0}^{T-1} \rho_{t+1}
        \EE{}{\ell(\bw_t) - \inf \ell}
        \leq 
        \frac{\rho_0}{2 \eta \sum_{t=0}^{T-1} \rho_{t+1}}
        \EE{}{\Vert \bw_0 - \bw_* \Vert^2}
        + \eta \frac{\sigma^2 + (\delta + G)^2}{2} 
        +  \varepsilon^{-1} \frac{\delta^2}{2}.
    \end{equation*}
    Now define $\bar \bw_T : =\frac{1}{\sum_{t=0}^{T-1} \rho_{t+1}}\sum_{t=0}^{T-1} \rho_{t+1} \bw_t = \frac{1}{\sum_{t=0}^{T-1} \theta^t} \sum_{t=0}^{T-1} \theta^t \bw_t$, and use Jensen's inequality to get
    \begin{equation*}
        \EE{}{\ell(\bar \bw_T) - \inf \ell}
        \leq 
        \frac{\Vert \bw_0 - \bw_* \Vert^2}{2 \eta \sum_{t=0}^{T-1} \rho_{t+1}}
        + \eta \frac{\sigma^2 + (\delta + G)^2}{2} 
        +  \varepsilon^{-1} \frac{\delta^2}{2}.
    \end{equation*}
    We will now simplify this upper bound, so we can later make an appropriate choice of  $\eta$ leading to the desired complexity bound.
    We will now focus on the geometric sum appearing in the denominator, and lower bound it.
    Start by using the definition of $\rho_t$ to write
    \begin{equation*}
        \eta \sum_{t=0}^{T-1} \rho_{t+1}
        =
        \eta \theta \sum_{t=0}^{T-1} \theta^{t}
        =
        \eta \frac{\theta}{1- \theta} (1 - \theta^T).
    \end{equation*}
    On the one hand, we see that from the definition of $\theta$ we have
    \begin{equation*}
        \eta \frac{\theta}{1- \theta}
        =
        \frac{\eta }{1 + \eta \varepsilon} \frac{1}{1 - \frac{1}{1 + \eta \varepsilon}}
        =
        \frac{\eta }{\eta \varepsilon}
        =
        \varepsilon^{-1}.
    \end{equation*}
    On the other hand, we can guarantee that $(1 - \theta^T) \geq 1/2$ provided that $\varepsilon =  \frac{2}{\eta T}$.
    To see this, first observe that $(1 - \theta^T) \geq 1/2$ is equivalent to $T \geq \frac{\ln(2)}{\ln(\theta^{-1})}$.
    Since we impose that $\eta \varepsilon = \frac{2}{T}$ then necessarily  $\theta = \frac{T}{T+2}$ which means that $\theta^{-1} = 1 + \frac{2}{T}$.
    So what we need to verify is $T \geq \frac{\ln(2)}{\ln(1 + \frac{2}{T})}$.
    Now, observe that $\ln(2) \leq 1$.
    Moreover, it is an exercise (see \cref{L:log inequality 1/2} which is deferred to the end of this proof) to verify that for all $x \geq 0$, $\frac{1}{\ln(x+1   )} \leq \frac{1}{2} + \frac{1}{x}$.
    So for our condition to hold it is enough to ask that $T \geq \frac{1}{2} + \frac{T}{2}$, which is equivalent to $T\geq 1$, which is true.
    
    Combining all those bounds together with our new definition for $\varepsilon$, we now have
    \begin{equation*}
        \EE{}{\ell(\bar \bw_T) - \inf \ell}
        \leq 
        \frac{2 \Vert \bw_0 - \bw_* \Vert^2}{\eta T}
        + \eta \frac{\sigma^2 + (\delta + G)^2}{2} 
        +  \frac{\eta \delta^2 T}{4}.
    \end{equation*}
    
    To simplify the last stage of our analysis, let us note  $a= 2\Vert \bw_0 - \bw_* \Vert^2$, $b=\frac{\sigma^2 + (\delta + G)^2}{2}$ and $c=\frac{\delta^2}{4}$, so that our  bound writes as
    \begin{equation}\label{cvsmgdc2}
    \EE{}{\ell(\bar \bw_T) - \inf \ell}
        \leq 
        \frac{a}{\eta T} + \eta b  + c \eta T.
    \end{equation}
    Minimizing the right-hand side with respect to $\eta$ is equivalent to solve $\frac{a}{\eta T} = \eta b  + c \eta T$, where
    \begin{equation*}
        \frac{a}{\eta T} = \eta b  + c \eta T
        \Leftrightarrow
        \eta^2 = \frac{a}{T(b+cT)}
        \Leftrightarrow
        \eta = \sqrt{\frac{a}{bT+cT^2}}
        \Leftrightarrow
        \eta = 2\sqrt{\frac{2\Vert \bw_0 - \bw_* \Vert^2}{2(\sigma^2 + (\delta + G)^2)T+\delta^2 T^2 }}
    \end{equation*}
    With such choice of stepsize, our bound becomes
    \begin{equation*}
    \EE{}{\ell(\bar \bw_T) - \inf \ell}
        \leq 
        \frac{2a}{\eta T}
        =
        \frac{2a}{T} \sqrt{\frac{bT+cT^2}{a}}
        =
        2\sqrt{a}\sqrt{\frac{b}{T}+c}
        =
        2\Vert \bw_0 - \bw_* \Vert \sqrt{\frac{\sigma^2 + (\delta + G)^2}{T}+\frac{\delta^2}{2}}
    \end{equation*}
    and this gives us our first bound.
    For the second bound, let us take simply $\eta = \frac{1}{\delta T + \sqrt{T}}$ and inject it into \eqref{cvsmgdc2} to obtain
    \begin{eqnarray*}
    \EE{}{\ell(\bar \bw_T) - \inf \ell}
        & \leq  &
        \frac{a \delta T + a \sqrt{T}}{ T} + \frac{b}{\delta T + \sqrt{T}}  +  \frac{cT}{\delta T + \sqrt{T}} \\
        & \leq &
        a \delta +
        \frac{a}{ \sqrt{T}} 
        + \frac{b}{\sqrt{T}}  +  \frac{c}{\delta}
        =
        \frac{a + b}{\sqrt{T}} + a \delta + \frac{c}{\delta} \\
        & \leq  &
        \frac{2D^2 + \sigma^2 + \delta^2 + G^2}{\sqrt{T}} + (2D^2 +1)\delta, 
    \end{eqnarray*}
    where in the last inequality we simplified some numerical constants.
\end{proof}

\begin{lemma}\label{L:log inequality 1/2}
    For every $x \geq 0$, $\frac{1}{\ln(1+x)} \leq \frac{1}{2} + \frac{1}{x}$.
\end{lemma}

\begin{proof}
    This inequality is equivalent to $\ln(1+x) \geq \frac{2x}{x+2}$, or again $(x+2)\ln(1+x) \geq 2x$.
    Define $\varphi : (-1,+\infty) \to \mathbb{R}$ as $\varphi(x) = (x+2)\ln(1+x)$ and compute its derivatives:
    \begin{equation*}
        \varphi'(x) = \ln(1+x) +1+\frac{1}{1+x} 
        \quad \text{ and } \quad 
        \varphi''(x) = \frac{x}{(1+x)^2}.
    \end{equation*}
    We see that $\varphi''(x) \geq 0$ for all $x \geq 0$, so $\varphi$ is convex on $[0,+\infty)$.
    So we can use the tangent inequality:
    \begin{equation}
        \varphi(x) \geq \varphi(0) + \varphi'(0)(x-0) = 0 + 2x = 2x.
    \end{equation}
    Therefore $\varphi(x) \geq 2x$ for all $x \geq 0$, which is what we wanted to prove.
\end{proof}




\subsection{Rates for \SMGD{}  for Convex Lipschitz Functions: Proof of \cref{P:CV SMGD convex lipschitz} }



\CVSMGDconvexlipschitz*

\begin{proof}
    In this proof, we note $\mathcal{F}_t$ the filtration generated by $\bw_0, \dots, \bw_t$, and we will note $\mathbb{E}_t$ to refer to the expectation conditioned on $\mathcal{F}_t$.
    
    We want to apply \cref{T:CV SMGD for convex lipschitz}.
    Let us denote $\bg_t = \m_t + \be_t$, so that the iterates of \SMGD{} verify $\bw_{t+1} = \bw_t - \eta \bg_t$.
    We are now going to show that $\bg_t$ has uniformly bounded bias and variance.
    For this we will make use of \cref{ass:uniform-bias-variance-bound-median,Ass:errors bounded expectation}.
    Remember that we assumed $\delta_2=\sigma_2=0$ without loss of generality (we can do this according to \cref{L:oracle purely bouded if lipschitz objective}), so we note $\delta$ and $\sigma$ instead of $\delta_1$ and $\sigma_1$.
    Also, \cref{Ass:errors bounded expectation} on the errors together with the tower rule imply that $\EE{t}{\Vert \be_t \Vert^2} = \EE{t}{\E{\Vert \be_t \Vert^2 \ | \ \bg_t^{(1)}, \dots, \bg_t^{(n)} }}  \leq \varepsilon^2$. 
    Regarding the bias, we can write 
    \begin{equation*}
        \Vert \EE{t}{\bg_t - \nabla \ell(\bw_t)} \Vert 
        \leq
        \Vert \EE{t}{\m_t - \nabla \ell(\bw_t) }\Vert 
        +
        \Vert \EE{t}{\be_t} \Vert 
        \leq
        \delta
        + \varepsilon,
    \end{equation*}
    where in the last inequality we used Jensen's inequality to write $\Vert \EE{t}{ \be_t }\Vert \leq \sqrt{\EE{t}{ \Vert \be_t \Vert^2}}$.
    As for the variance, we write
    \begin{equation*}
        \EE{t}{\Vert \bg_t - \mathbb{E}_{t} \bg_t \Vert^2}
        \leq 2\EE{t}{\Vert \m_t - \EE{t}{\m_t} \Vert^2}
        + 2 \EE{t}{\Vert \be_t - \EE{t}{\be_t} \Vert^2}
        \leq 
        2 \sigma^2 + 8 \EE{t}{\Vert \be_t \Vert^2} 
        \leq 
        2 \sigma^2 + 8 \varepsilon^2.
    \end{equation*}
    Using \cref{T:CV SMGD for convex lipschitz}, we obtain
    \begin{equation*}
        \EE{}{\ell(\bar \bw_T) - \inf \ell}
        \leq 
        \frac{2 \Vert \bw_0 - \bw_* \Vert^2}{\eta T}
        + \eta \frac{2\sigma^2 + 8\varepsilon^2 + (\delta + \varepsilon + G)^2}{2} 
        +  \frac{\eta (\delta + \varepsilon)^2 T}{4}.
    \end{equation*}
    Setting the step size $\eta = \frac{1}{(\delta + \varepsilon) T + \sqrt{T}}$, we conclude that 
    \begin{equation*}
        \EE{}{\ell(\bar \bw_T) - \inf \ell}
        \leq 
        \frac{2D^2 + 2\sigma^2 + 8 \varepsilon^2+ (\delta + \varepsilon)^2 + G^2}{\sqrt{T}} + (2D^2 +1)(\delta + \varepsilon).
    \end{equation*}
\end{proof}


\subsection{Rates for \SMGD{} with Stable Noise: Proof of \cref{T:CV SMGD stable noise}}

%
\corstableshort*

\begin{proof}
First of all, remember that the $\ell_1$-median can be computed exactly, by computing component-wise a one-dimensional median.
This allows us to have $\varepsilon =0$ in \cref{Ass:errors bounded expectation}. Further, due to \cref{stable noise implies finite bias variance v2} we have $\delta_1=0$ in \cref{ass:uniform-bias-variance-bound-median}.

\textbf{Proof of part (i).} Choose $\eta = \frac{1}{8L(1+2\sigma_2+2\delta_2)\sqrt{T}}$, where $\delta_2,\sigma_2$ are from \cref{stable noise implies finite bias variance v2}. Hence, we can apply \cref{prop:approx-median-descent}. As $\eta = \frac{C}{\sqrt{T}}$ with $C=\frac{1}{8L(1+2\sigma_2+2\delta_2)}$, we get the bound $\frac{1}{C\sqrt{T}} + \frac{C}{\sqrt{T}} = \mathcal{O}(1/\sqrt{T})$.

\textbf{Proof of part (ii).}  The result is a direct consequence of \cref{stable noise implies finite bias variance v2} and \cref{P:CV SMGD convex lipschitz}.

This completes the proof.
\end{proof}



\section{Approximate Computation of Gradient Estimators}\label{sec:app-more-update-formulas}
%
\subsection{Gradient Estimation Problem: General Considerations}
%
Given a function $\dist: \R^d \to \R$ we consider the problem \eqref{prob:gradient-learning}, which writes as 
\begin{equation*}
   \argmin{\vec{m} \in \R^d} \E{ \dist(\vec{m}-\vec{z})}.
\end{equation*}
In what follows, we will mostly consider functions $\dist$ satisfying the following assumption.

\begin{assumption}\label{Ass:D convex has solutions}
    The function $\dist : \mathbb{R}^d \to \mathbb{R}$ is convex, $\argmin{\vec{z}} \mathcal{D}(\vec{z})= \{\vec{0}\}$ and problem \eqref{prob:gradient-learning} admits a solution.
\end{assumption}

Note that assuming $\mathcal{D}$ to be convex and finite implies that it is continuous, therefore measurable, which means that $\E{ \dist(\vec{m}-\vec{z})}$ is well defined.
The assumption that $\argmin{\vec{z}}\mathcal{D}(\vec{z})= \{\vec{0}\}$ guarantees that in the trivial case when the distribution is a Dirac $\delta_{\m_0}$, then the solution to  \eqref{prob:gradient-learning} is exactly $\m_0$.
The assumption that a solution to \eqref{prob:gradient-learning} exists is easily verified in practice, under the assumption that $\dist$ is coercive:

\begin{lemma}\label{L:D coercive implies median exists}
    Let $\dist : \mathbb{R}^d \to \mathbb{R}$ be convex.
    Assume that $\mathbb{E}\left[ \mathcal{D}(\bz - \mathbb{E}\left[ \bz \right]) \right] < + \infty$. 
    If $\dist$ is coercive, i.e. $\lim\limits_{\|\vec{m}\| \to \infty} \dist(\vec{m})= + \infty$, then \eqref{prob:gradient-learning} admits a solution.
\end{lemma}

\begin{proof}
    Let $\varphi(\m):= \mathbb{E}\left[ \mathcal{D}(\m - \bz) \right]$.
    Since we assume that $\varphi(\mathbb{E}\left[ \bz \right]) = \mathbb{E}\left[ \mathcal{D}(\bz - \mathbb{E}\left[ \bz \right]) \right] < + \infty$ then $\varphi$ is proper.
    Since $\mathcal{D}$ is convex then $\varphi$ is convex, as an expectation of convex functions.
    % For the lower semicontinuity this is a consequence of Fatou's lemma.
    Moreover $\mathcal{D}$ is convex and finite, which means that $\mathcal{D}$ is continuous on $\mathbb{R}^d$, see e.g. Corollary 8.39 in \citet{BauCom}.
    Therefore $\varphi$ is an expectation of  lower semicontinuous functions, which implies that $\varphi$ is lower semicontinuous.
    
    Since $\mathcal{D}$ is coercive, convex and continuous, we know that there exists $a \in (0,+\infty)$ and $b \in \mathbb{R}$ (see Proposition 14.16 in \citet{BauCom}) such that 
    $\mathcal{D}(\cdot) \geq a \Vert \cdot \Vert + b$.
    Therefore
    \begin{eqnarray*}
        \varphi(\m) & = & \mathbb{E}\left[ \mathcal{D}(\m - \bz) \right] 
        \geq
        a \mathbb{E}\left[ \Vert \m - \bz \Vert  \right] + b \\
        & \geq &
        a \Vert \m - \mathbb{E}\left[ \bz \right] \Vert - a\mathbb{E}\left[ \Vert \bz - \mathbb{E}\left[ \bz \right] \Vert \right] + b \\
        & \geq &
        a \Vert \m - \mathbb{E}\left[ \bz \right] \Vert - a\mathbb{E}\left[ \mathcal{D}( \bz - \mathbb{E}\left[ \bz \right] ) \right]
        \underset{\Vert \m \Vert \to \infty}{\longrightarrow} + \infty,
    \end{eqnarray*}
    where in the inequalities we used the fact that $\Vert \cdot \Vert$ is $1$-Lipschitz, and that $a \mathbb{E}\left[ \Vert \bz - \mathbb{E}\left[ \bz \right] \Vert \right] \leq a \mathbb{E}\left[ \mathcal{D}(\bz - \mathbb{E}\left[ \bz \right]) \right] -b < + \infty$.
    So $\varphi$ is coercive on top of being proper convex and lower semicontinuous. We can then conclude that $\varphi$ has a minimizer, with for instance Proposition 11.15 from \citet{BauCom}.
\end{proof}

%% OLD MODE PARAGRRAPH
% \paragraph{Mean, Median, Mode.} As indicated in \cref{sec:spp-updates}, if $\mathcal{D} = \tfrac12 \Vert \cdot \Vert^2$, then the solution to \eqref{prob:gradient-learning} is the \emph{mean}, if instead $\mathcal{D}$ is a $\ell_p$-norm, then the solution to \eqref{prob:gradient-learning} is the $\ell_p$-\emph{median} of the distribution.
% We also show here that the \emph{mode} of a distribution can be seen as a solution of \eqref{prob:gradient-learning}, when $\mathcal{D}$ is a particular nonconvex function.

% \begin{lemma}
%     Let $\mathcal{D}$ be the $\ell_0$-pseudonorm $\Vert \cdot \Vert_0$, where $\Vert \m \Vert_0$ returns the number of nonzero coefficients of $\m$.
%     Assume that the distribution of $\bz$ is independent component-wise.
%     Then the component-wise \emph{modes} of $\bz$ is a solution to \eqref{prob:gradient-learning}, which are the values appearing with the largest probability.
% \end{lemma}

% \begin{proof}
%     First, note that $\Vert \cdot \Vert_0$ is separable, in the sense that $\Vert \m \Vert_0 = \sum_{j=1}^d \vert \m_i \vert_0$, where $\vert t \vert_0 = 0$ if $t=0$ and $1$ otherwise.
%     So \eqref{prob:gradient-learning} can be decomposed component-wise, which means that we only have to prove the result for $d=1$.
%     And $\mathbb{E}\left[ \vert m - z \vert_0 \right]$ is exactly $\mathbb{P}(z \neq m)$, so minimizing it is equivalent to  find $m$ which maximizes $\mathbb{P}(z=m)$.
% \end{proof}


\subsection{Approximations Based on the Stochastic Proximal Point: General $\mathcal{D}$}\label{SS:SPP:general D formulas}

Given a proper, closed convex function $\varphi : \mathbb{R}^d \to \mathbb{R}\cup\{+\infty\}$ and $\tau > 0$, we recall the definition of the proximal operator
\begin{equation*}
    \prox{\tau \varphi}(\m_0) = \argmin{\m \in \mathbb{R}^d} \varphi(\m) + \frac{1}{2 \tau} \Vert \m - \m_0 \Vert^2.
\end{equation*}
Applying the \SPP{} algorithm to \eqref{prob:gradient-learning-3} gives the following iterations 
\begin{equation}\label{D:SPP median D abstract}
    \begin{cases}
        \text{Sample $\bg_t$ i.i.d. from $\mathcal{G}$} \\
    \m_{t+1} = \argmin{\m \in \mathbb{R}^d} \dist(\vec{m} - \vec{g}_t) + \frac{1}{2\tau}\|\vec{m}-\vec{m}_t\|^2.
    %\prox{\tau \mathcal{D}(\cdot - \bg_t)}(\m_t).
    \end{cases}
\end{equation}
As we see next, those iterations can be reformulated to simply involve the proximal operator of $\mathcal{D}$:

\begin{restatable}{proposition}{viewpointsSPP}\label{prop:viewpoints}
    Let $\mathcal{D}$ verify \cref{Ass:D convex has solutions}.
    The \SPP{} update~\eqref{D:SPP median D abstract} for solving~\eqref{eqn:online-spp-update} can be equivalently written as
    \begin{align}
        \vec{m}_{t+1} &=  \vec{g}_t + \prox{\tau\dist}(\vec{m}_t-\vec{g}_t) \label{eq:gtprox} \\
        &=  \prox{\tau\dist^*(\frac{\vec{m}_t - \cdot }{\tau})}(\vec{g}_t) , \label{eq:mtprox}
    \end{align}
    where $\dist^*(\vec{m}) := \sup_{\vec{y}} \big( \dotprod{\vec{y},\vec{m}} -  \dist(\vec{y}) \big)$ is the Fenchel conjugate of $\dist$, and $\bg_t$ is sampled i.i.d. from $\mathcal{G}$ at each iteration.
 \end{restatable}
 \begin{proof}
 Here we make use of proximal calculus rules, which can found in most textbooks on convex analysis, for example Chapter 6 in~\citet{Beck2017}.
 Starting from \eqref{D:SPP median D abstract}, and applying the variable transformation $\by=\bar \by+\vec{g}_t$, we have 
 \begin{align*} 
     \begin{split}
         &\vec{m}_{t+1} = \argmin{\by\in \R^d} \dist(\by-\vec{g}_t) + \tfrac{1}{2\tau} \|\by-\vec{m}_t\|^2 \\
         &\hspace{1ex}=\vec{g}_t + \argmin{\bar \by\in \R^d} \dist(\bar \by) + \tfrac{1}{2\tau} \|\bar \by-(\vec{m}_t-\vec{g}_t)\|^2  \\
        &\hspace{1ex}=\vec{g}_t + \prox{\tau\dist}(\vec{m}_t-\vec{g}_t).
     \end{split}
 \end{align*}
 This proves~\eqref{eq:gtprox}. As for~\eqref{eq:mtprox}, we combine $(\tau \dist)^*=\tau\dist^*(\cdot/\tau)$ \citep[Thm.\ 4.14]{Beck2017} and the Moreau decomposition theorem \citep[Thm.\ 6.44]{Beck2017}, to obtain
 \[ \prox{\tau\dist}(\vec{m}) = \vec{m} - \prox{\tau\dist^*(\frac{\cdot}{\tau})}(\vec{m}). \]
 Plugging into~\eqref{eq:gtprox} we have
     \begin{align}
        \vec{m}_{t+1} &=  \vec{g}_t + \prox{\tau\dist}(\vec{m}_t-\vec{g}_t) \nonumber \\
       &=  \vec{g}_t + \vec{m}_t-\vec{g}_t - \prox{\tau\dist^*(\frac{\cdot}{\tau})}(\vec{m}_t-\vec{g}_t)  \nonumber \\
       &= \vec{m}_t- \prox{\tau\dist^*(\frac{\cdot}{\tau})}(\vec{m}_t-\vec{g}_t)   \label{eq:templz8oehdzd}
        \end{align}
       
        Finally, using that for any function $f$ it holds $\argmin{\vec{y}} f(\vec{y}) = \vec{m} + \argmin{\bar{\vec{y}}} f(\bar{\vec{y}}+\vec{m})$
        in~\eqref{eq:templz8oehdzd} gives
       \begin{align*}
        \vec{m}_{t+1} &=   \vec{m}_t- \left(\vec{m}_t+\prox{\tau\dist^*(\frac{\cdot+\vec{m}_t}{\tau})}(-\vec{g}_t)  \right)\\
        &= -\prox{\tau\dist^*(\frac{\cdot+\vec{m}_t}{\tau})}(-\vec{g}_t) \\
        &= \prox{\tau\dist^*(\frac{\vec{m}_t\cdot -}{\tau})}(\vec{g}_t) 
        \end{align*}
         where in the final equality we used that if $f(\vec{x}) = g(-\vec{x})$ we have that
        \[\prox{g}(\vec{m}) = - \prox{ f }(-\vec{m}). \]
\end{proof}


\subsection{Approximations Based on the Stochastic Proximal Point: Particular Cases of $\mathcal{D}$}
\label{app:sec:SPP-updates-particular-cases}


We first characterize the \SPP{} update~\eqref{eqn:online-spp-update} for when $\dist$ is the $\ell_p$-norm, and show it simply require to project the sampled gradient onto a certain ball of radius $\tau$ centered at $\vec{m}_t$. 

\begin{restatable}{proposition}{SPPLp}\label{prop:pnormview}
    Let $\dist = \|\cdot\|_p$ be the $\ell_p$-norm for $p \in [1, \infty].$ Let $q \in [1, \infty]$ such that $\frac{1}{p} + \frac{1}{q}=1$, and let 
    $\mathbb{B}_{q}(\vec{m},\tau)  \; := \; \left\{ \vec{y} \; : \; \norm{\vec{m}-\vec{y}}_q \leq \tau \right\}.$
    The \SPP{} update~\eqref{eqn:online-spp-update} is given by
    \begin{align}
        \vec{m}_{t+1} &=  \vec{m}_t \; +\; \Proj{\mathbb{B}_{q}(\vec{0},\tau)  } \big( \vec{g}_t-\vec{m}_t \big) \label{eq:projzero} \\
        &= \Proj{\mathbb{B}_q( \vec{m}_t,\tau) }\big(  \vec{g}_t \big).\label{eq:projmt}
    \end{align}
\end{restatable}
%
\begin{proof}
The proximal operator of the $\ell_p$-norm is given by 
$$\prox{\tau \norm{\cdot}_p}(\vec{m})  = \vec{m} - \tau \Proj{\mathbb{B}_{q}(\vec{0}) }(\vec{m}/\tau),$$
see Example 6.47 in~\citet{Beck2017}. Using this in~\eqref{eq:gtprox} gives
\begin{align*}
 \vec{m}_{t+1} &=    \vec{g}_t  + \prox{\tau\norm{\cdot}_p}(\vec{m}_t-\vec{g}_t) \\
 &=  \vec{m}_t  -\tau \Proj{\mathbb{B}_{q}(\vec{0},1) }\big(\frac{\vec{m}_t-\vec{g}_t}{\tau}\big)\\
  &=  \vec{m}_t  +\tau \Proj{\mathbb{B}_{q}(\vec{0},1) }\big(\frac{\vec{g}_t-\vec{m}_t}{\tau}\big)\\
    &=  \vec{m}_t  + \Proj{\mathbb{B}_{q}(\vec{0},\tau) }\big( \vec{g}_t-\vec{m}_t \big)\\
  &=  \vec{m}_t  +\left( \Proj{\mathbb{B}_{q}(\vec{m}_t,\tau) }\big( \vec{g}_t \big)- \vec{m}_t\right)\\
  % &= \tau \Proj{\mathbb{B}_{q}(\frac{\vec{m}_t}{\tau},1) }\big(\frac{ \vec{g}_t}{\tau}\big) \\
  &= \Proj{\mathbb{B}_{q}( \vec{m}_t,\tau) }\big(  \vec{g}_t \big) 
\end{align*}
where the third equality we used that 
$$\Proj{\mathbb{B}_{q}(\vec{0},1)(\vec{m})} = -\Proj{\mathbb{B}_{q}(\vec{0},1)(-\vec{m})},  $$ followed by
\[\Proj{\mathbb{B}_{q}(\vec{m}/\tau,1)}(\vec{y}/\tau) = \frac{1}{\tau} \Proj{\mathbb{B}_{q}(\vec{m},\tau)}(\vec{y})   \]
in the  fourth equality and 
\[  \Proj{\mathbb{B}_{q}(\vec{0},\tau)}(\vec{y}) = \Proj{\mathbb{B}_{q}(\vec{m},\tau)}(\vec{y}-\vec{m}) - \vec{m}  \]
in the fifth equality.
\end{proof}
%
%
%
\begin{figure}
    \centering
    \includegraphics[width=7cm,, clip]{figures/proj_new_tau.png}
    \caption{One \SPP{} update~\eqref{eq:projmt} for $\dist=\|\cdot\|_1$ (blue) or $\dist=\|\cdot\|_2$ (green) amounts to project onto a ball of the dual norm, centered at the current iterate $\vec{m}_t$ and of radius $\tau > 0$. The thick lines represent the projections of $\bg_t$ for all possible values of $\tau$. }
    \label{fig:proj}
\end{figure}

Using \cref{prop:pnormview} we can now develop closed form updates of the \SPP{} method for when $\dist$ is the $\ell_1$- or the $\ell_2$-norm, with both methods being related to clipping.


%%%%%%%%%%%%%%%%%%
\corLonenorm*
%%%%%%%%%%%%%%%%%
\begin{proof}
        Follows from~\eqref{eq:projzero} and that
        \[ \;\min\{\max\{\vec{v}, - \tau\} , \tau\}=  \Proj{\mathbb{B}_{\infty}( \vec{0},\tau) }\big(  \vec{v} \big). \]
        % that is $\vec{m}_{t+1} =  \vec{g}_t + \prox{\tau\|\cdot\|_2}(\vec{m}_t - \vec{g}_t)$, and \cref{lem:prox-and-clip-l2}. \qed \\ 
\end{proof}
%
%


%%%%%%%%%%%%%%%%%%%%%%%%%
\corLtwonorm*
%%%%%%%%%%%%%%%%%%%%%%%%%
   \begin{proof}
        Follows from~\eqref{eq:projzero},  by using that 
        \[\Proj{\mathbb{B}_{2}( \vec{0},\tau) }\big(  \vec{v} \big) = \frac{\tau \; \vec{v}}{\max\{\tau, \|\vec{v}\|_2 \}} . \]
        %A direct proof is also given in \cref{lem:prox-and-clip-l1}. 
    \end{proof}


We end this section with the closed form updates of the \SPP{} method for when $\dist$ is the Huber function.

%%%%%%%%%%%%%%%%%%%%%%%%%
\corLhuber*
%%%%%%%%%%%%%%%%%%%%%%%%%

\begin{proof}
    Follows from \eqref{eqn:online-spp-update} and the fact that the proximal operator of the Huber function is given by 
    $\prox{\tau H_{\mu}}(\vec{z}) = \big(1-\frac{\mu\tau}{\max\{\|\vec{z}\|, \mu(1+\tau)\}}\big) \vec{z}.$ For a proof, see \citet[Example 6.66]{Beck2017}, and note that here we have the additional factor $\mu$ in the definition of $H_{\mu}$.
\end{proof}


%



\subsection{Approximations Based on Subgradient Updates}
\label{sec:SGDmedian}
Instead of \SPP{}, we could solve~\eqref{prob:gradient-learning-3} by iteratively taking steps of stochastic subgradient descent. If $\dist$ is convex, then in each iteration we sample a gradient $\vec{g}_t$, and update our current estimate $\vec{m}_t$ by
\begin{align}\label{eqn:subgradient-step}
    \vec{m}_{t+1} = \vec{m}_t - \tau\vec{u}_t,\quad \vec{u}_t \in \partial \dist(\vec{m}_t - \vec{g}_t).
\end{align}
Here, $\tau >0$ is the learning rate, and $\vec{u}_t$ is a subgradient (of the convex subdifferential). Since $\dist$ could be a non-differentiable function such as the $\ell_2$- or $\ell_1$-norm, we use subgradients instead of gradients. 

\paragraph{Squared $\ell_2$-norm.} As a first simple example of~\eqref{eqn:subgradient-step}, let $\dist = \tfrac{1}{2}\norm{\cdot}_2^2.$ In this case~\eqref{eqn:subgradient-step} becomes
\begin{align} \label{eq:mom1}
    \vec{m}_{t+1} = \vec{m}_t - \tau(\vec{m}_t - \vec{g}_t) = (1-\tau) \vec{m}_t + \tau\vec{g}_t.
\end{align}
This is again (heavy-ball) momentum (cf.\ \ref{eqn:momentum}), now with coefficient $\beta=1-\tau$. Note that in contrast to \SPP{}, the coefficient $\beta$ could be negative, and is only in $[0,1)$ (as typically is the case for momentum) if $\tau \in [0,1]$.  

\paragraph{$\ell_1$-norm.} If we choose $\dist= \|\cdot\|_1$, update \eqref{eqn:subgradient-step} gives
\begin{align*}
    \vec{m}_{t+1} = \vec{m}_t + \tau\sgn(\vec{g}_t - \vec{m}_t),
\end{align*}
where the $\sgn$-operator can take any value in $[-1,1]$ for coordinates $i\in [d]$ such that $(\vec{g}_{t}-\vec{m}_{t})_i=0$.
%
This recovers a version of sign-\SGD{} \citep{Riedmiller1993,pmlr-v80-bernstein18a}, where the $\mathrm{sign}$-operation is applied to the increment $\vec{m}_t - \vec{g}_t$. For example, if we reset $\vec{m}_t=0$ in every iteration, we obtain exactly sign-\SGD{}.

\paragraph{$\ell_2$-norm.} If we choose $\dist= \|\cdot\|_2$, and if $\vec{g}_t \neq \vec{m}_t$, update \eqref{eqn:subgradient-step} gives
    \begin{align*}
        \vec{m}_{t+1} = \vec{m}_t - 
            \tau\frac{\vec{m}_t-\vec{g}_t}{\|\vec{m}_t-\vec{g}_t\|} = \big(1-\frac{\tau}{\|\vec{m}_t-\vec{g}_t\|}\big)\vec{m}_t + \frac{\tau}{\|\vec{m}_t-\vec{g}_t\|} \vec{g}_t.
    \end{align*}
    Otherwise, if $\vec{g}_t = \vec{m}_t$, we can choose any $\vec{u}_t$ with $\|\vec{u}_t\|\leq 1$ and set $\vec{m}_{t+1} = \vec{m}_t - \tau \vec{u}_t$; in particular we can set $\vec{m}_{t+1} = \vec{m}_t$.


\subsection{Rates for \SPP{}}

Now that we have presented how to compute the iterates \eqref{eqn:online-spp-update} in practice, let us state its convergence properties. Since all these methods are instantiates of the \SPP{} method, we can use standard results such as \citet[Thm.\ 4.4]{Davis2019}.

\begin{restatable}{proposition}{proxformedian}
\label{prop:prox for computing median}
    Let $\dist$ be a norm, let $\varphi(\vec{m}) := \EE{\vec{g}\sim \mathcal{G}}{ \dist(\vec{m}- \vec{g})}$, and let $\vec{m}_* \in \argmin{\vec{m}\in \R^d} \varphi(\vec{m})$. 
    If $\vec{m}_t$ is generated by \eqref{eqn:online-spp-update}
    then for $\bar{\vec{m}}_T = \frac{1}{T}\sum_{t=1}^T \vec{m}_t$ we have
    \begin{equation*}
        \mathbb{E}\left[ \varphi(\bar{\vec{m}}_T) - \varphi(\vec{m}_*) \right] = O \left( \frac{\Vert \vec{m}_0 - \vec{m}_* \Vert^2}{\tau T} + \tau \right).
    \end{equation*}
\end{restatable}
%
\begin{proof}
    Since we assume that $\dist$ is a norm over $\R^d$ then it convex and continuous. 
    Moreover, since all norms are equivalent in $\mathbb{R}^d$, we have that $\dist(\vec{w}) \leq C \norm{\vec{w}}_2$ for some constant $G>0$. Thus $\dist$ is a $G$-Lipschitz function. 
    If now we write $\varphi_{\bg} := \mathcal{D}( \cdot - \bg)$, it is clear that it is convex and Lipschitz.
    Therefore, $\varphi = \EE{\bg}{\varphi_{\bg}}$ is itself a convex $G$-Lipschitz function.
    The claim now follows from \citet[Theorem 4.4]{Davis2019}.
\end{proof}

Here is a corollary stating rates on the expected distance between the approximate sample median and the true sample median, when using the $\ell_1$-norm:

\begin{corollary}
    In the context of \cref{prop:prox for computing median}, assume that $\mathcal{D}= \|\cdot\|_1$ and that $\mathcal{G}$ is a sum of Dirac measures $\tfrac{1}{n}\sum_{i=1}^n \delta_{\bg^{(i)}}$.
    Then we have further
    \begin{equation*}
        \mathbb{E}\left[ \Vert \bar{\vec{m}}_T - \vec{m}_* \Vert \right] = O \left( \frac{\Vert \vec{m}_0 - \vec{m}_* \Vert^2}{\tau T} + \tau \right).
    \end{equation*}
\end{corollary}


\begin{proof}
    In this proof, we use the notation $(\vec{x}_j)_{j=1,\ldots,d}$ to specify the indivdual components of a vector $\vec{x}\in \R^d$.
    For $\vec{x} \in \mathbb{R}^d$, we define the sign operator as follows: let $\sgn(\vec{x}) = (\sgn(\vec{x}_j))_{j=1,\ldots,d}$, where $\vec{x}_j$ is the $j$-th coordinate of $\vec{x} \in \R^d$, and where we define for a scalar $t\in \R$ the set-valued operator $\sgn(t) = \{+1\}$ if $t>0$, $\sgn(t) = \{-1\}$ if $t<0$, and $\sgn(t) = [-1,1]$ if $t=0$.
    Further, given a set $A \subset \mathbb{R}^d$ we define $\Vert A \Vert_- := \inf\limits_{ a \in A} \Vert a \Vert$.

    To obtain the desired bound, and given that we already have the conclusion of \cref{prop:prox for computing median}, it is enough for us to show that there exists $\mu >0$ such that 
    \begin{equation*}
        (\forall \m \in \mathbb{R}^d) \quad \mu \Dist(\m, {\rm{argmin}}~\varphi) \leq \varphi(\m) - \inf \varphi.
    \end{equation*} 
    It is a standard result from variational analysis that this property\footnote{It is sometimes referred to as  (weak) sharp minima \citep{Fer91,BurFer93}, (superlinear)(linear) conditioning \citep{Lem92,CorJouZal97,Lem98}, or error bound \citep{LewPan98}.} is true if and only if the following inequality is satisfied (see Theorem 5.1 in \citet{CorJouZal97} or Proposition 3.1 in \citet{Lem98})
    \begin{equation*}
         (\forall \m \in \dom \partial \varphi) \quad \Vert \partial \varphi(\m) \Vert_- \geq \mu.
    \end{equation*}
    Our goal now is to prove the above inequality, from which the conclusion will follow. 
    Let us consider $\m \notin {\rm{argmin}}~\varphi$, and use standard calculus to write 
    \begin{equation*}
        \partial \varphi(\m) = \frac{1}{n} \sum_{i=1}^n \sgn(\m-{\bg}_i)
        =
        \frac{1}{n}
        \sum_{i=1}^n
        %
        \big(\sgn(\m_j-({\bg}_i)_j)\big)_{j=1,\ldots,d}
        =
        \frac{1}{n}
        %  
        \Big(\sum_{i=1}^n \sgn(\m_j-({\bg}_i)_j) \Big)_{j=1,\ldots,d}.
    \end{equation*}
    Our goal is to show that $\Vert \partial \varphi(\m) \Vert_- \geq \frac{1}{n}$, which is equivalent to show that $\Vert\sum_{i=1}^n \sgn(\m-{\bg}_i) \Vert_- \geq 1$. 
    Because $\m \notin {\rm{argmin}}~\varphi$, we know that $0 \notin \partial \varphi(\m)$ so there must exist a coordinate $j$ such that the $j$-th coordinate of $\partial \varphi(\m)$ does not contain zero, in other words such that $0 \notin \sum_{i=1}^n \sgn(\m_j-({\bg}_i)_j)$.
    Using the fact that $\Vert \cdot \Vert_2 \geq \Vert \cdot \Vert_\infty$, we can therefore lower bound $\Vert\sum_{i=1}^n \sgn(\m-{\bg}_i) \Vert_- \geq \vert \sum_{i=1}^n \sgn(\m_j-({\bg}_i)_j) \vert_-$.
    Let us denote $s_i := \sgn(\m_j - ({\bg}_i)_j)$, so that we want to prove that $\vert \sum_{i=1}^n s_i \vert_- \geq 1$.
    Let us now consider a few cases.
    \begin{itemize}
        \item If $s_i = \{+1\}$ or $\{-1\}$ for every $i$, then $\sum_{i=1}^n s_i$ is a singleton. Furthermore it is a sum of relative numbers, so $\sum_{i=1}^n s_i \in \mathbb{Z}$. But we also know that $0 \notin \sum_{i=1}^n s_i$, so we conclude that $\sum_{i=1}^n s_i \in \mathbb{Z}^*$, and so that $\vert \sum_{i=1}^n s_i \vert_- \geq 1$.
        \item If $s_i = [-1,+1]$ for every $i$, then we immediately see that $0 \in \sum_{i=1}^n s_i$ which is a contradiction.
        \item Otherwise, the $s_i$ are combinations of singletons and intervals. Let us note $I = \{i : s_i = [-1,+1]\}$ and $I' = \{i : s_i = \{ \pm 1\}\}$, which are not empty by assumption.
        Let us also note $k \geq 1$ the cardinality of $I$, so that $\sum_{i \in I} s_i = [- k , +k]$.
        We also introduce $s := \sum_{i \in I'} s_i \in \mathbb{Z}$.
        Now we can write
        \begin{equation*}
            \vert \sum_{i=1}^n s_i \vert_-
            =
            \inf\limits_{ t \in \sum_{i=1}^n s_i} \vert t \vert 
            =
            \inf\limits_{t \in [-k,+k]} \vert t + s \vert
            =
            d(s;[-k,+k]),
        \end{equation*}
        where the latter is the distance from the number $s$ to the interval $[-k,+k]$.
        Now remember that $0 \notin  \sum_{i=1}^n s_i$, so this distance cannot be zero, which means that $\vert s \vert > k$.
        Moreover both $s$ and $k$ are rational numbers, so  this distance must be greater or equal to $1$. This concludes the proof.
    \end{itemize}
\end{proof}






%%%%%%%%%%%%%%%%%%%%%%%%%%%%%%%%%%%%%%%%%%%%%%
\section{Supplementary Information on Experiments}\label{sec:app-experiments-supp}
This section provides additional information for the experimental setup.

For the language modeling experiments, all details are identical to \citet{kunstner2023noise}, Section A.1. They also provide implementations for all tasks at \url{https://github.com/fKunstner/noise-sgd-adam-sign}, which we use.
For the training runs for language modeling, we use batch size $256$ for \texttt{PTB}, $320$ for \texttt{WikiText-2}, and $32$ for \texttt{SQuAD}.

\subsection{Additional Plots}\label{sec:app-fixed-weight-transformer}

\paragraph{Ablations for \cref{fig:toy-example}}. We provide two ablation studies for the experiment on gradient estimation on heavy-tailed synthetic data. First, \cref{fig:toy-example-ablation} (left) shows that for Gaussian noise ($\alpha=2$), the performance of momentum can surpass \texttt{V/CClip} depending on selection of the step size $\tau$. However, this does not fix the issue of a quickly degrading performance for momentum, when $\alpha$ decreases.

Second, \cref{fig:toy-example-ablation} (right) shows the performance under \emph{skewed} heavy-tailed noise; for this, we set the skewness parameter of the $\alpha$-stable distribution to one.

\begin{figure}[H]
    \centering
    \begin{subfigure}{0.37\columnwidth}  
        \includegraphics[width=0.99\textwidth]{figures/toy_example/levy_ablation_all_noise.pdf}
    \end{subfigure}
    \hspace{2cm}
    \begin{subfigure}{0.37\columnwidth}  
        \includegraphics[width=0.99\textwidth]{figures/toy_example/levy_skewed_all_noise.pdf}
    \end{subfigure}
    \caption{\textbf{(Left)} Same as \cref{fig:toy-example} but with $\tau=0.001$. Here, momentum performs best for Gaussian data, but still becomes instable when the noise is more heavy-tailed. \textbf{(Right)} Same as \cref{fig:toy-example} but with \emph{skewed} noise. Here, the median of the noise is no longer zero, which explains the failure of \texttt{V/CClip}.}
    \label{fig:toy-example-ablation}
\end{figure}
%
%
%

\begin{figure}[H]
    \centering
    \includegraphics[width=0.99\textwidth]{{figures/least_squares/huber_ablation/huber_alpha_1.1_independent_False_state_dependent_False_scale_1.0}.pdf}
    \caption{Least-squares problem (same as \ref{item:independent}) for \HC{} with different values of $\mu$ and $\tau$. Setting $\tau=1.0$ performs relatively well across the selected range of values of $\mu$.}
    \label{fig:huber-mu-ablation}
\end{figure}


\paragraph{Full gradient estimation with fixed weights.} Here, we show the results for estimating the full gradient of a transformer architecture, with weights fixed at initialization.
\cref{fig:ptb-fixed-weights} shows the results for \texttt{PTB} where we tried three different batch sizes $\{64,256,1024\}$. \cref{fig:wikitext-fixed-weights} shows the results for \texttt{WikiText-2}.

We remark that in this experiment only, we have used the first order approximation of \eqref{eqn:momentum-as-spp}, namely $\vec{m}_{t+1} =  \left(1-\tau\right) \vec{m}_t + \tau \vec{g}_t$. That is, we replace $\frac{\tau}{1+\tau}$ with its first-order Taylor approximation around zero, which is equal to $\tau$. As the value of $\tau$ is small, this has negligible impact on the result: for example, if $\tau=0.01$, then $\frac{\tau}{1+\tau} \approx 0.0099$.

\emph{Discussion.} From \cref{fig:ptb-fixed-weights}, we observe two phenomena: in the long run, momentum attains the lowest error for estimating the full batch gradient. However, the initial decrease of the error is much faster for \CC{}, followed by \VC{}. This is important when using these estimates 
within a training setup such as \eqref{eq:online-update},
where we only do one iteration of gradient estimation, followed by an update of the weight (and hence a change in the full-batch gradient). Secondly, we observe that the difference in convergence speeds is most pronounced when the batch size increases. Hence, being robust to outliers seems not to be solvable only by increasing the batch size and thus decreasing the noise of the mini-batch gradient. In fact, the contrary seems to be the case. This is similar to the observations made in \citet{kunstner2023noise}.

\begin{figure}[H]
    \centering
    \includegraphics[width=0.99\columnwidth]{figures/transformer_example/fixed_weights/ptb_transformer_encoder.pdf}
    \caption{Encoder transformer on \texttt{PTB} dataset, with weights fixed at initialization. We use momentum, $\cclip{\tau}$ and $\vclip{\tau}$ to estimate the full-batch gradient.}
    \label{fig:ptb-fixed-weights}
\end{figure}

\begin{figure}[H]
    \centering
    \includegraphics[width=0.8\textwidth]{figures/transformer_example/fixed_weights/wikitext2_transformer_xl.pdf}
    \caption{Encoder transformer on \texttt{Wikitext-2} dataset, with weights fixed at initialization. We use momentum, $\cclip{\tau}$ and $\vclip{\tau}$ to estimate the full-batch gradient.}
    \label{fig:wikitext-fixed-weights}
\end{figure}

\subsection{Learning Rate Values for Transformer Training}

Here, we report for each of the language modelling tasks, the learning rate value that is displayed in \cref{fig:transformer-training}.
For each method, we tuned the learning rate over a set of values $10^{\frac{j}{2}}$ for a suitable interval of integer numbers $j \in \mathbb{Z}$.

% changed LR of SGDM and clipped-SGD to adjust for missing dampening in code
\begin{table}[h]
    \caption{Tuned learning rate values used for each method. All methods use constant learning rates. For \SGDM{}, tuning information is also reported in Section C.1 in \citep{kunstner2023noise}).}
    \vspace{0.5cm}
    \centering
    \begin{tabular}{c|ccccc}
         Name&  \SGDM{} & \texttt{clipped-SGD} & \VC{} & \CC{} & \texttt{Adam} \\
         \hline \\
         \texttt{PTB} & 1.0 & 3.16 & 3.16 & 0.316 & 0.001 \\
         \texttt{WikiText-2}& 0.316 & 3.16 & 1.0 & 0.316 & 0.001 \\
         \texttt{SQuAD}& 0.316 & 1.0 & 0.316 & 0.1 & 0.000316\\
    \end{tabular}
    \label{tab:transformer-lr-values}
\end{table}

%%%%%%%%%%%%%%%%%%%%%%%%%%%%%%%%%%%%%%%%%%%%%%%%%%%%%%%%%%%%%%%%%%%%

\end{document}
